\documentclass[11pt,a4paper]{article}
\usepackage[utf8]{inputenc}
\usepackage[T1]{fontenc}
\usepackage{lmodern}
\usepackage{amsmath,amssymb,amsthm}
\usepackage{geometry}
\usepackage{hyperref}
\usepackage{fancyhdr}
\usepackage{mathrsfs}
\usepackage{array}
\usepackage{booktabs}
\usepackage{longtable}
\usepackage{graphicx}

\geometry{margin=1in}

\pagestyle{fancy}
\fancyhf{}
\rhead{Cayley --- Collected Papers, Vol.\ V}
\lhead{Pages 551--600}
\cfoot{\thepage}

\title{\textbf{Mathematical Papers}\\\Large Arthur Cayley --- Collected Papers, Volume V\\\normalsize Pages 551--600}
\author{}
\date{}

\newtheorem*{theorem}{Theorem}
\newtheorem*{lemma}{Lemma}
\renewcommand{\between}{\mathbin{|}}

\begin{document}

\maketitle
\thispagestyle{empty}

\tableofcontents
\newpage


\section*{379. Notices of Communications to the British Association for the Advancement of Science.}

[From the Reports of the British Association for the Advancement of Science.]

\subsection*{5. On Curves of the Third Order. Report, 1861, p.~2.}

A curve of the third order or cubic curve is a section of a cubic cone and such
cone is intersected by a concentric sphere in a spherical cubic. It is an obvious
consequence of a theorem of Sir Isaac Newton's that there are five principal kinds
of cubic cones, or what is the same thing five principal kinds of spherical cubics---
but the nature of these five kinds of spherical cubics was first distinctly explained by
M\"{o}bius. They may be designated the \textit{simplex}, the \textit{complex}, the \textit{crunodal}, the \textit{acnodal}
and the \textit{cuspidal}: where crunode, acnode, denote respectively the two species of double
points (nodes), viz.\ the double point with two real branches, and the conjugate or
isolated point. The foregoing results are known: the special object of the paper is
to establish a subdivision of the simplex kind of spherical cubics. The simplex kind
is a continuous reentering curve cutting a great circle, to fix the ideas say the
equator, in three pairs of opposite points, which are the three real inflexions of the
curve. The three great circles which are the tangents at the inflexions and the
equator divide the entire surface of the sphere into fourteen regions whereof eight
are trilateral and the remaining six are quadrilateral. The curve may be entirely in
six out of the eight trilateral regions, and it is in this case said to be \textit{simplex\noindent trilateral}; or it may lie entirely in the six quadrilateral regions, and it is in this
case said to be \textit{simplex quadrilateral}; and there is an intermediate form, the \textit{simplex neutral}; viz.\ in this case the three great circles tangents at the inflexions meet in
a pair of opposite points and there are in all only twelve regions all of them
trilateral; the curve lies entirely in six of these regions.

\subsection*{6. On a Certain Curve of the Fourth Order. Report, 1862, p.~3.}

The curve in question is the locus of the centres of the conics which pass
through three given points and touch a given line; if the equations of the sides of
the triangle formed by the three points are $x=0$, $y=0$, $z=0$, these coordinates being
such that $x+y+z=0$ is the equation of the line infinity, and if $\alpha x+\beta y+\gamma z=0$ be
the equation of the given line, then (as is known) the equation of the curve is
\[
\sqrt{\alpha x\,(y+z-x)} + \sqrt{\beta y\,(z+x-y)} + \sqrt{\gamma z\,(x+y-z)} = 0.
\]
The special object of the communication was to exhibit the form of the curve in
the case where the line cuts the triangle, and to point out the correspondence of the
positions of the centre upon the curve, and the point of contact on the given line.


\subsection*{7. On the Representation of a Curve in Space by means of a Cone and Monoid Surface. Report, 1862, p.~3.}

The author gave a short account of his researches recently published in the
\textit{Comptes Rendus}. The difficulty as to the representation of a curve in space is, that
such a curve is not in general the complete intersection of two surfaces; any two
surfaces passing through the curve intersect not only in the curve itself, but in a
certain companion curve, which cannot be got rid of; this companion curve is in the
proposed mode of representation reduced to the simplest form, viz.\ that of a system
of lines passing through one and the same point. The two surfaces employed for the
representation of a curve of the $n$th order are, a \textit{cone} of the $n$th order having for
its vertex an arbitrary point (say the point $x=0$, $y=0$, $z=0$), and a \textit{monoid} surface
with the same vertex, viz.\ a surface the equation whereof is of the form $Qw-P=0$,
$P$ and $Q$ being homogeneous functions of $(x, y, z)$ of the degrees $p$ and $p-1$
respectively (where $p$ is at most $=n-1$). The monoid surface contains upon it
$p\,(p-1)$ lines given by the equations $(P=0,\ Q=0)$; and, the cone passing through
$n\,(p-1)$ of these lines (if, as above supposed, $p=n-1$, this implies that some of
these lines are multiple lines of the cone), the monoid surface will besides intersect
the cone in a curve of the $n$th order.

\subsection*{8. On a Formula of M.~Chasles relating to the Contact of Conics. Report, 1864, p.~1.}

The author gave an account of the recent investigations of M.~Chasles in relation
to the theory of conics, viz., M.~Chasles has found that the properties of a system
of conics, containing one arbitrary parameter, depend upon two quantities called by
him the \textit{characteristics} of the system; these are, $\mu$, the number of conics of the
system which pass through a given point, and, $\nu$, the number of conics of the system
which touch a given line; or, say, $\mu$ is the \textit{parametric order}, $\nu$ the \textit{parametric class},
of the system. And he exhibited a transformation obtained by him of a formula of
M.~Chasles for the number of conics which touch five given curves, viz., if $(M, m)$,
$(N, n)$, $(P, p)$, $(Q, q)$, $(R, r)$ be the orders and classes of the five given curves respec-\noindent tively, then the number of curves is
\[
=(1,\ 2,\ 4,\ 4,\ 2,\ 1)\,(M,\ m)\,(N,\ n)\,(P,\ p)\,(Q,\ q)\,(R,\ r),
\]
where the notation stands for $1.MNPQR + 2\Sigma mNPQR + 4\Sigma mnPQR + \&$c. The trans-
formed formula in question was communicated by the author to M.~Chasles, and had
appeared in the \textit{Comptes Rendus}; but it is, in fact, included in a very beautiful and
general theorem given in the same Number by M.~Chasles himself.

\subsection*{9. On the Problem of the In-and-circumscribed Triangle. Report, 1864, p.~1.}

The general problem of the in-and-circumscribed triangle may be thus stated, viz.,
to find a triangle the angles whereof severally lie in, and the sides severally touch, a
given curve or curves; and we may, in the first instance, inquire as to the number
of such triangles. The first and easiest case is when the curves are all distinct; here,
if the angles lie in curves of the \textit{orders} $m$, $n$, $p$, respectively, and the sides touch
curves of the \textit{classes} $Q$, $R$, $S$, respectively, then the number of triangles is $=2mnpQRS$.
The number may be obtained for some other cases; but the author has not yet
considered the final and most difficult case, viz.\ that in which the angles severally
lie in, and the sides severally touch, one and the same given curve.

\medskip

The foregoing notices relate to verbal communications upon questions with which
I was at the time occupied and which are for the most part more fully discussed in
papers printed elsewhere. I remark upon them as follows:

\begin{enumerate}
\item[1.] I have no remembrance as to this; I think no paper printed or written.

\item[2.] See 175.

\item[3.] See 158.

\item[4.] No paper printed. The intention was to consider the different forms of
trinodal quartic curves, in particular those with real nodes, as obtained from the
inversion of a conic according to the different relations of the conic to the fundamental
triangle. Thus according as the conic cuts in two real points, touches, or cuts in two
imaginary points, a side of the triangle, the tangents at the corresponding node are
real, coincident, or imaginary; viz.\ the node is a crunode, cusp, or acnode. And in
the case of real intersections there is a further distinction according as the inter-
sections lie each or either of them on the side itself, or on the side produced in
one or other of the two directions. By considering the different relations of the conic
to the fundamental triangle we thus obtain the different forms of the trinodal quartic.

\item[5.] See 351.

\item[6.] I think no paper printed or written.

\item[7.] See 302 and 305.

\item[8.] See 306.

\item[9.] The question is considered in a memoir ``On the Problem of the In-and-circum-
scribed Triangle,'' \textit{Phil.\ Trans.}\ t.\ CLXI. (for 1871), pp.\ 369---412.
\end{enumerate}

\newpage


\section*{380. Note on the Rectangular Hyperbola.}

[From the \textit{Oxford, Cambridge and Dublin Messenger of Mathematics}, vol.\ \textsc{i}. (1862), p.\ 77.]

Every conic which passes through the points of intersection of two rectangular
hyperbolas is a rectangular hyperbola. In fact if a conic be referred to rectangular
axes, the condition that it may be a rectangular hyperbola is $\text{Coeff.\ of }x^2 = -\text{Coeff.\ of }y^2$.
Hence if $U$, $V$ be any two quadratic functions of $x$, $y$, and if $\lambda$ be a constant, the
condition in question being satisfied for each of the functions $U$, $V$, is satisfied for
the function $U+\lambda V$: and the equation of any conic through the points of intersection
of the conics $U=0$, $V=0$ is $U+\lambda V=0$: which proves the theorem in question.

In particular if from two of the angles of a triangle perpendiculars are let fall
on the opposite sides, and if the point of intersection of the perpendiculars and the
third angle be joined: then since the first side and the perpendicular upon it are a
rectangular hyperbola, and the second side and the perpendicular upon it are a
rectangular hyperbola; the third side and the joining line must be a rectangular
hyperbola: that is, these two lines must be at right angles to each other. We have
thus the well-known theorem that the perpendiculars let fall from the angles of a
triangle on the opposite sides meet in a point.

The theorem as to the hyperbolas is a particular case of the theorem that three
conics which pass through the same four points are met by any line whatever in
six points forming a system in involution. In fact a rectangular hyperbola is a conic
meeting the line at infinity in two points harmonically related to the circular points
at infinity: hence two of the conics being rectangular hyperbolas, the foci of the
involution are the circular points at infinity: hence these points and the points in
which the line at infinity meets the third conic are harmonically related to each other:
that is, the third conic is a rectangular hyperbola.

\newpage


\section*{381. Note on Bezout's Method of Elimination.}

[From the \textit{Oxford, Cambridge and Dublin Messenger of Mathematics}, vol.\ \textsc{ii}. (1864), pp.\ 88, 89.]

Let $U$, $U'$ be any two rational and integral functions of $x$ of the same order;
to fix the ideas let them be the cubic functions
\begin{align*}
U &= ax^3 + bx^2 + cx + d,\\
U' &= a'x^3 + b'x^2 + c'x + d'.
\end{align*}

Write
\[
\begin{array}{c@{\quad}c}
A = \begin{vmatrix} U, & U' \\ a, & a' \end{vmatrix}, &
P = \begin{vmatrix} U, & U' \\ a, & a' \end{vmatrix}, \\[2ex]
B = \begin{vmatrix} U, & U' \\ b, & b' \end{vmatrix}, &
Q = \begin{vmatrix} U, & U' \\ ax+b, & a'x+b' \end{vmatrix}, \\[2ex]
C = \begin{vmatrix} U, & U' \\ c, & c' \end{vmatrix}, &
R = \begin{vmatrix} U, & U' \\ ax^2+bx+c, & a'x^2+b'x+c' \end{vmatrix}, \\[2ex]
D = \begin{vmatrix} U, & U' \\ d, & d' \end{vmatrix}, &
S = \begin{vmatrix} U, & U' \\ ax^3+bx^2+cx+d, & a'x^3+b'x^2+c'x+d' \end{vmatrix} = \begin{vmatrix} U, & U' \\ U, & U' \end{vmatrix} = 0,
\end{array}
\]
then we have
\begin{align*}
P &= A,\\
Q &= Ax + B,\\
R &= Ax^2 + Bx + C,\\
S &= Ax^3 + Bx^2 + Cx + D = 0,
\end{align*}
and thence
\begin{align*}
A &= P,\\
B &= Q - Px,\\
C &= R - Qx,\\
D &= S - Rx = -Rx.
\end{align*}

Let $\alpha$ be an arbitrary quantity and write
\[
\square x =
\begin{vmatrix}
U & U' \\
a\alpha^3+b\alpha^2+c\alpha+d' & a'\alpha^3+b'\alpha^2+c'\alpha+d'
\end{vmatrix};
\]
we have it is clear
\begin{align*}
\square &= A\alpha^3 + B\alpha^2 + C\alpha + D,\\
&= \alpha^2 P + \alpha^2(Q-Px) + \alpha(R-Qx) + Rx,\\
&= (\alpha^3 - \alpha^2 x)\,P + (\alpha^2 - \alpha x)\,Q + (\alpha - x)\,R,
\end{align*}
and hence
\[
\frac{\square}{\alpha - x} = \alpha^2 P + \alpha Q + R.
\]

The equations $P=0$, $Q=0$, $R=0$ are respectively quadratic equations in $x$, the
equations which are used in Bezout's method of elimination; and representing them by
\begin{align*}
P &= Lx^2 + Mx + N = 0,\\
Q &= L'x^2 + M'x + N' = 0,\\
R &= L''x^2 + M''x + N'' = 0,
\end{align*}
we have
\[
\begin{vmatrix}
L, & M, & N \\
L', & M', & N' \\
L'', & M'', & N''
\end{vmatrix} = 0
\]
as the equation resulting from the elimination of $x$ from the equations $U=0$, $U'=0$.
The foregoing investigation shows that the functions $P$, $Q$, $R$ are obtained as the
coefficients of $\alpha^2$, $\alpha$, $1$ in the development of
\[
\frac{1}{\alpha - x}
\begin{vmatrix}
U & U' \\
a\alpha^3+b\alpha^2+c\alpha+d, & a'\alpha^3+b'\alpha^2+c'\alpha+d'
\end{vmatrix},
\]
or more generally, taking $U$, $U'$ to be any two functions of the order $n$, that the $n$
functions $P$, $Q$, $R$, \&c.\ each of the order $n-1$ are obtained as the coefficients of
$\alpha^{n-1}$, $\alpha^{n-2}$, \ldots $\alpha$, $1$ in the development of
\[
\frac{1}{\alpha - x}
\begin{vmatrix}
U, & U' \\
U_{\alpha}, & U'_{\alpha}
\end{vmatrix},
\]
where $U_{\alpha}$, $U'_{\alpha}$ are what $U$, $U'$ become when $x$ is replaced therein by $\alpha$: and we
have thus a simple \textit{\`{a} posteriori} verification of the form in which, several years ago,
I presented Bezout's Method of Elimination.

\begin{flushright}
2, \textit{Stone Buildings, W.C., March 5, 1863.}
\end{flushright}

\newpage


\section*{382. Note on the Tetrahedron.}

[From the \textit{Oxford, Cambridge and Dublin Messenger of Mathematics}, t.\ \textsc{iii}. (1866), pp.\ 8---10.]

The following simple properties of a tetrahedron seem worth noticing.

In the tetrahedron $ABCD$ if $AC=BD$ and $AD=BC$, then the line joining the
middle points of $AB$ and $CD$, or say the points $\frac12AB$ and $\frac12CD$, cuts at right angles
these lines $AB$ and $CD$.

If $AB=CD$, then the line joining the points $\frac12AC$, $\frac12BD$, and the line joining
the points $\frac12AD$, $\frac12BC$ (lines which in any tetrahedron meet each other), cut each
other at right angles.

In fact if $A$, $B$, $C$, $D$ have for their coordinates $(\alpha_1, \beta_1, \gamma_1)$, $(\alpha_2, \beta_2, \gamma_2)$, $(\alpha_3, \beta_3, \gamma_3)$,
$(\alpha_4, \beta_4, \gamma_4)$: then the coordinates of the point $\frac12AB$ are $\frac12(\alpha_1+\alpha_2)$, $\frac12(\beta_1+\beta_2)$, and so
for the points $\frac12CD$, \&c.; the equations of the line through the points $\frac12AB$, $\frac12CD$
therefore are
\[
\frac{x - \frac12(\alpha_1+\alpha_2)}{\alpha_1+\alpha_2-\alpha_3-\alpha_4}
= \frac{y - \frac12(\beta_1+\beta_2)}{\beta_1+\beta_2-\beta_3-\beta_4}
= \frac{z - \frac12(\gamma_1+\gamma_2)}{\gamma_1+\gamma_2-\gamma_3-\gamma_4},
\]
and I observe in passing that this line passes through the point whose coordinates are
\[
\tfrac14(\alpha_1+\alpha_2+\alpha_3+\alpha_4),\quad
\tfrac14(\beta_1+\beta_2+\beta_3+\beta_4),\quad
\tfrac14(\gamma_1+\gamma_2+\gamma_3+\gamma_4);
\]
the other two similar lines pass through the same point, and the above-mentioned
property of the general tetrahedron is thus proved.

The condition that the foregoing line may cut at right angles the line $AB$, the
equations whereof are
\[
\frac{x-\alpha_1}{\alpha_1-\alpha_2} = \frac{y-\beta_1}{\beta_1-\beta_2} = \frac{z-\gamma_1}{\gamma_1-\gamma_2},
\]
is at once seen to be
\[
\Sigma\,(\alpha_1-\alpha_2)(\alpha_1+\alpha_2-\alpha_3-\alpha_4) = 0,
\]
where $\Sigma$ denotes the sum of the corresponding terms in $\alpha$, $\beta$, $\gamma$. And so the
condition that the same line may cut at right angles the line $CD$ is
\[
\Sigma\,(\alpha_3-\alpha_4)(\alpha_1+\alpha_2-\alpha_3-\alpha_4) = 0.
\]

But the conditions $AC=BD$, $AD=BC$ give respectively
\[
\Sigma\,[(\alpha_1-\alpha_3)^2 - (\alpha_2-\alpha_4)^2] = 0,\quad
\Sigma\,[(\alpha_1-\alpha_4)^2 - (\alpha_2-\alpha_3)^2] = 0,
\]
or writing these in the form
\begin{align*}
\Sigma\,(\alpha_1-\alpha_3-\alpha_2+\alpha_4)(\alpha_1+\alpha_3-\alpha_2-\alpha_4) &= 0,\\
\Sigma\,(\alpha_1-\alpha_4-\alpha_2+\alpha_3)(\alpha_1+\alpha_4-\alpha_2-\alpha_3) &= 0,
\end{align*}
we obtain, by successively adding and subtracting, the two required equations.

The equations of the line through $\frac12AC$, $\frac12BD$ are
\[
\frac{x-\frac12(\alpha_1+\alpha_3)}{\alpha_1+\alpha_3-\alpha_2-\alpha_4}
= \frac{y-\frac12(\beta_1+\beta_3)}{\beta_1+\beta_3-\beta_2-\beta_4}
= \frac{z-\frac12(\gamma_1+\gamma_3)}{\gamma_1+\gamma_3-\gamma_2-\gamma_4},
\]
and those of the line through $\frac12AD$, $\frac12BC$ are
\[
\frac{x-\frac12(\alpha_1+\alpha_4)}{\alpha_2+\alpha_4-\alpha_1-\alpha_3}
= \frac{y-\frac12(\beta_1+\beta_4)}{\beta_2+\beta_4-\beta_1-\beta_3}
= \frac{z-\frac12(\gamma_1+\gamma_4)}{\gamma_2+\gamma_4-\gamma_1-\gamma_3};
\]
and the condition that these may cut at right angles is
\[
\Sigma\,(\alpha_1+\alpha_4-\alpha_2-\alpha_3)(\alpha_1+\alpha_4-\alpha_2-\alpha_3) = 0,
\]
that is
\[
\Sigma\,\{(\alpha_1-\alpha_2)^2 - (\alpha_3-\alpha_4)^2\} = 0,
\]
which is in fact the condition $AB=CD$.

Combining the two theorems we see that if in a tetrahedron the pairs of opposite
sides are respectively equal, then the line joining the centres of opposite sides cuts
these sides at right angles, and moreover the three joining lines cut each other at
right angles.

A tetrahedron of the form in question may be constructed as follows: viz.\ taking
a parallelogram $ABCD$, whereof the diagonals $AC$, $BD$ are unequal, then bending the
parallelogram about its shorter diagonal $AC$ in such manner that in the solid figure
$BD$ becomes equal to $AC$, we have a tetrahedron the opposite sides whereof are
respectively equal.

Or it may be constructed even more simply as follows: viz.\ if $ABCD'$ and
$A'B'C'D$ be parallel faces of any rectangular parallelopiped (the angles $A$ and $A'$,
$B$ and $B'$, $C$ and $C'$, $D$ and $D'$ being respectively opposite to each other), then $ABCD$
or $A'B'C'D'$ is a tetrahedron of the form in question. The consideration of the
rectangular parallelopiped puts in evidence the foregoing geometrical property.

In such a tetrahedron the line joining the centres of a pair of opposite sides
is in the language of Bravais, see his ``M\'{e}moire sur les poly\`{e}dres de forme sym\'{e}trique,''
\textit{Liouville}, t.\ \textsc{xiv}. (1849), pp.\ 141---180, a binary axis of symmetry: viz.\ the figure is
not altered by turning it round such axis through an angle $=\frac12\cdot 360^\circ$. There are thus
three such axes at right angles to each other, but the figure has not any centre of
symmetry, nor (assuming that it is not further particularised) any plane of symmetry:
each of the three axes is a principal axis, and the figure belongs to the sixth of
Bravais' twenty-three classes of polyhedra, see the table p.\ 179. It was in fact by
seeking to construct a figure of this class that I was led to the investigation.

\newpage


\section*{383. Problems and Solutions.}

[From the \textit{Mathematical Questions with their Solutions from the Educational Times}, vols.\ \textsc{i.}\ to \textsc{iv.}, 1863 to 1865.]

\subsection*{[Vol.~\textsc{i}. (June 1863 to June 1864), pp.\ 18, 19.]}

\paragraph*{1373.} (By T.~T.~Wilkinson, F.R.A.S.)---Given a circle $(C)$ and any point $A$, either
within or without the circle: through $A$ draw $BAD$ cutting the circle in $B$, $D$.
Then it is required to find another point $E$, such that, if $LEM$ be drawn cutting
the circle in $L$, $M$, we may always have $AE^2 = LE \cdot EM \pm BA \cdot AD$.

\begin{center}
\textit{Solution by Professor Cayley.}
\end{center}

Consider a circle centre $O$ and radius $OA$, and in relation thereto a point $M$
either outside or inside the circle, and suppose that
\[
(OA)^2 - (OM)^2, \text{ or the ``squared inner potency'' of $M$ is denoted by } \square i \cdot M,
\]
and
\[
(OM)^2 - (OA)^2, \text{ or the ``squared outer potency'' of $M$ is denoted by } \square o \cdot M,
\]
so that, for an outside point, $\square o \cdot M = -\square i \cdot M$, is the square of the tangential
distance of $M$ from the circle; and, for an inside point, $\square i \cdot M = -\square o \cdot M$, is the
square of the shortest semi-chord through $M$.

Suppose now that $M$ is a given point; the proposed question is in effect to find
the locus of a point $P$ such that $\pm\, \square o \cdot P \pm \square o \cdot M = (MP)^2$; but we have thus in
reality four different questions according as the signs are assumed to be $++$, $+-$, $-+$
or $--$; the case $++$, or when $\square o \cdot P + \square o \cdot M = (MP)^2$, is perhaps the most inter-
\noindent esting.

Taking the radius as unity, $(\alpha, \beta)$ as the coordinates of $M$, and $(x, y)$ as the
coordinates of $P$, we have here
\[
(x^2+y^2-1) + (\alpha^2+\beta^2-1) = (x-\alpha)^2 + (y-\beta)^2, \text{ or } \alpha x + \beta y - 1 = 0;
\]
that is, the locus of $P$ is a right line, the polar of $M$ in regard to the circle.

It may be remarked, that, when $M$ is an inside point, then throughout the locus
$P$ is an outside point; and, replacing the negative quantity $\square o \cdot M$ by its value,
$=-\square i \cdot M$, we have $\square o \cdot P - \square i \cdot M = (MP)^2$. If, however, $M$ is an outside point, then
in part of the locus $P$ is an outside point, and we have $\square o \cdot P + \square o \cdot M = (MP)^2$, while
in the remainder of the locus $P$ is an inside point, and, replacing the negative
quantity $\square o \cdot P$ by its value, $=-\square i \cdot P$, we have $-\square i \cdot P + \square o \cdot M = (MP)^2$. For the
case $+-$, the locus of $P$ is a right line, but for each of the other two cases
$-+$ and $--$ the locus is a circle; the discussion of the several cases presents no
particular difficulty.


\subsection*{[Vol.~\textsc{i}. pp.\ 43---45.]}

\paragraph*{1387.} (By W.~K.~Clifford.)---1. Four common tangents are drawn to a circle
and an ellipse which passes through the centre $(O)$ of the circle; if $A$, $B$ be opposite
intersections of the tangents, prove that $OA$ and $OB$ are equally inclined to the
tangent at $O$ to the ellipse.

2. If a straight line $A$ join the poles of $B$ with respect to two conics, prove
that the lines joining $AB$ to a pair of opposite intersections of common tangents,
form, with $A$, $B$, an harmonic pencil.

3. If a point $A$ be the intersection of the polars of $B$ with respect to two
conics, and $AB$ be cut by a pair of common chords in $C$, $D$, prove that $ACBD$ is
an harmonic range.

\begin{center}
\textit{2. Solution by Professor Cayley.}
\end{center}

This elegant theorem is included as a particular case in the known theorem,
``Given three conics inscribed in the same quadrilateral, the tangents from any point
to these conics form a pencil in involution.''

Mr Clifford's theorem is in fact as follows: viz., Four common tangents are drawn
to a circle and an ellipse which passes through the centre $O$ of the circle; if $A$, $B$
be opposite intersections of the tangents, then $OA$, $OB$ are equally inclined to the
tangent at $O$ to the ellipse.

This comes to saying that the tangent at $O$ to the ellipse, say $OT$, is the double
or sibi-conjugate line of the involution of the pencil formed by the lines $OA$, $OB$,
and the lines $OI$, $OJ$ drawn from $O$ to the circular points at infinity; and if we
replace the circle by an arbitrary conic $S$, and the line at infinity by an arbitrary
line $IJ$, the theorem will be as follows:

Consider a conic $S$; a line meeting this conic in the points $I$, $J$; and the point
$O$, the intersection of the tangents at $I$, $J$, or (what is the same thing) the pole of
the line $IJ$ in regard to the conic. If through the point $O$ there be drawn any
other conic $\Theta$, and if $A$, $B$ be opposite intersections of the common tangents of the
conics $S$, $\Theta$; then the tangent $OT$ at the point $O$ to the conic $\Theta$ is the double or
sibi-conjugate line of the involution of the pencil formed by the lines $OA$, $OB$, and
the lines $OI$, $OJ$; or, as we may also express it, the lines $OT$, $OT$, the lines $OA$, $OB$,
and the lines $OI$, $OJ$ form a pencil in involution.

Now, considering the two points or point-pair $(A, B)$ as a conic inscribed in the
quadrilateral formed by the common tangents of the conics $S$ and $\Theta$, the conics
$S$ and $\Theta$ and the point-pair $(A, B)$ are a system of three conics inscribed in the
same quadrilateral; and hence, by the general theorem above referred to, if $O'$ be any
point whatever, the tangents from $O'$ to the conic $S$, the tangents from $O'$ to the
conic $\Theta$, and the tangents from $O'$ to the point-pair (that is, the two lines
$O'A$, $O'B$) form a pencil in involution. But, if $O'$ coincide with $O$, then the tangents
to the conic $S$ are the lines $OI$, $OJ$; and the tangents to the conic $\Theta$ are the
coincident lines $OT$, $OT$; and we have thence the theorem in question; viz.\ that the
lines $OT$, $OT$, the lines $OI$, $OJ$, and the lines $OA$, $OB$ form a pencil in involution.

\subsection*{[Vol.~\textsc{i}. pp.\ 77---79.]}

\paragraph*{1409.} (By W.~K.~Clifford.)---For every point $A$ on a conic section there exists
a straight line $BC$, not meeting the curve, such that, if through any other point $K$
on the conic there be drawn any two straight lines meeting $BC$ in $B$, $C$, and the
curve in $D$, $E$, the angles $BAC$, $DAE$ are either equal or supplementary.

\begin{center}
\textit{Solution by Professor Cayley.}
\end{center}

I find that this very elegant theorem depends on the lemma to be presently
stated, and that it is intimately connected with Newton's theorem for the organic
description of a conic, or, what is the same thing, with the theorem of the anharmonic
relation of the points of a conic.

\begin{center}
\textsc{Lemma.}
\end{center}

If $AT$ be the tangent, and $AS$ any other line through a point $A$ of a
conic, and if two lines equally inclined to $AT$ and $AS$ respectively meet the conic
in the points $K$ and $D$ (viz., if $\angle TAK = SAD$, the two angles being measured in
opposite directions from $AT$, $AS$ respectively); then the line $KD$ meets $AS$ in a
fixed point $B$, that is, a point the position of which is independent of the magnitude
of the equal angles.

To prove this, take $A$ for the origin, and the bisectors of the angle $TAS$ for
the axes of $x$ and $y$: then the equation of the conic is
\[
ax^2 + 2hxy + by^2 + 2fy + 2gx = 0;
\]
the equation of the tangent at the origin, that is, the line $AT$, is $gx+fy=0$; and
hence the equation of the line $AS$ is $gx-fy=0$. Taking $y=\alpha x$ for the equation
of the line $AK$, we have, for the coordinates $x_1$, $y_1$ of the point $K$ where this meets
the conic,
\[
(a + 2h\alpha + b\alpha^2)\,x_1 + 2(f\alpha + g) = 0, \quad y_1 = \alpha x_1;
\]
and then the equation of the line $AD$ will be $y=-\alpha x$, and we shall have, for the
coordinates $x_2$, $y_2$ of the point $D$ where this meets the conic,
\[
(a - 2h\alpha + b\alpha^2)\,x_2 + 2(f\alpha + g) = 0, \quad y_2 = -\alpha x_2.
\]

The equation of the line $KD$ is
\[
\begin{vmatrix}
x, & y, & 1 \\
x_1, & \alpha x_1, & 1 \\
x_2, & -\alpha x_2, & 1
\end{vmatrix} = 0,
\]
that is
\[
\alpha x\,(x_1+x_2) + y\,(x_2-x_1) - 2\alpha x_1 x_2 = 0;
\]
and for the coordinates of the point $B$ where this meets the line $AS$, the equation
whereof is $gx-fy=0$, we have
\[
x\,\{f\alpha(x_1+x_2) + g(x_2-x_1)\} - 2f\alpha x_1 x_2 = 0,
\]
or, as this may be written,
\[
x\left\{\frac{f\alpha+g}{x_1} - \frac{-f\alpha+g}{x_2}\right\} - 2f\alpha = 0.
\]
But we have
\[
\frac{f\alpha+g}{x_1} = -\tfrac12(a+2h\alpha+b\alpha^2), \quad
\frac{-f\alpha+g}{x_2} = -\tfrac12(a-2h\alpha+b\alpha^2);
\]
and hence the equation is
\[
x(-2h\alpha) - 2f\alpha = 0,
\]
giving $x = -\dfrac{f}{h}$, and thence $y = -\dfrac{g}{h}$, for the coordinates of the point $B$; and, these
being independent of $\alpha$, the lemma is seen to be true.


Consider now the points $A$, $K$ as fixed points on the conic; then, revolving about
$A$ the constant angle $DAB$, and about $K$ the constant (zero) angle $DKB$, the locus
of $B$ is (by the theorem of the anharmonic relation of the points of a conic) given
in the first instance as a conic through the points $A$, $K$; but, observing that
a position of the angle $DAB$ is $TAK$, and that the corresponding position of $DKB$
is $AKA$, the line $AK$ is part of the locus; and the locus is made up of this line
and a line $BC$. And, conversely, given the fixed points $A$, $K$, and the line $BC$, the
original conic is, by Newton's theorem, described by means of the constant angles
$DAB$, $DKB$ revolving about these points in such a manner that the arms $AB$, $KB$
generate by their intersections the line $BC$. This being so, the other two arms
$AD$, $KD$ generate by their intersections the conic.

And then, considering the two positions $DAB$, $EAC$ of the angle $DAB$ (so that
$D$, $B$ are in a line with $K$, and $E$, $C$ are also in a line with $K$), we have
$\angle DAB = \angle EAC$, that is, $\angle DAE = \angle BAC$, which is Mr Clifford's theorem.

It has been seen that, $A$ being given, the same line $BC$ is obtained whatever
be the position of the point $K$; and, taking $AK$ for the normal at $A$, it at once
appears geometrically that (as remarked by Mr Clifford) the line $BC$ is the polar of
the point $\Theta$ of intersection of all the chords which subtend a right angle at $A$.

[Professor Cayley's lemma may be otherwise proved, as follows:

The trilinear equation of the conic, referred to two tangents ($\alpha$ at $A$, $\beta$ at $S$)
and their chord of contact ($\gamma$ or $AS$), is $U=\lambda\alpha\beta-\gamma^2=0$; and the equation of two
straight lines ($AK$, $AD$) equally inclined to $\alpha$, $\gamma$ is
\[
(\alpha - \mu\gamma)(\mu\alpha - \gamma) = 0, \quad \text{or}\quad V = \alpha^2 + \gamma^2 - (\mu + \mu^{-1})\,\alpha\gamma = 0;
\]
also $U+V=0$ denotes a conic passing through the intersections of $U$ and $V$; but
$U+V$ is resolvable into $\alpha=0$, or the tangent $AT$, and
\[
\alpha + \lambda\beta - (\mu + \mu^{-1})\,\gamma = 0,
\]
which is, therefore, the equation of the chord $KD$; whence we see that $KD$ meets
$AS$ (or $\gamma$) in a point $B$ (given by $\gamma=0$, $\alpha+\lambda\beta=0$) whose position is independent of
$\mu$, that is, of the equal angles $SAD$, $TAK$.]

\subsection*{[Vol.~\textsc{i}. pp.\ 125---127.]}

\paragraph*{1478.} (By J.~McDowell, M.A.)---$(\alpha)$ Two sides of a given triangle always pass
through two fixed points; prove that the third side always touches a fixed circle.

$(\beta)$ Two sides of a given triangle touch two fixed circles; prove that the third
side also touches a fixed circle.

$(\gamma)$ Two sides of a given polygon touch fixed circles; prove that all the remaining
sides also touch fixed circles.

\begin{center}
\textit{3. Solution by Professor Cayley.}
\end{center}

Since the theorem $(\gamma)$ follows at once from $(\beta)$, and $(\alpha)$ is included in $(\beta)$,
it is only necessary to prove $(\beta)$. Consider three given circles, and let it be proposed
to construct a triangle the sides whereof touch the given circles, and which is similar
to a given triangle; the direction of one side may be assumed at pleasure, and then
the triangle is determined. Impose now on the triangle the condition that the area
is equal to a given quantity; we obtain for the given area an expression involving
the angle $\theta$ which fixes the direction of one of the sides, and we have thus an
equation for the determination of the angle $\theta$. But, for a properly determined
relation between the data of the problem, the expression for the area becomes
independent of the angle $\theta$, that is, every triangle, the sides whereof touch the three
circles, and which is similar to a given triangle, is of the same area, or say, the area
of every such triangle is equal to a given quantity $\Delta$; and, this being so, it is
clear that, if we construct a triangle similar to a given triangle and of the given
area $\Delta$ (that is, a triangle equal to a given triangle), in such manner that two of
the sides touch two of the given circles, then the envelope of the remaining side will
be the remaining given circle; which is in fact the theorem $(\beta)$.

It only remains therefore to show that the foregoing porismatic case of the problem
exists.

For the first circle, let the coordinates of the centre be $a$, $b$, and the radius
be $c$; and suppose in like manner that we have $a'$, $b'$, and $c'$ for the second circle,
and $a''$, $b''$, and $c''$ for the third circle. Let $\lambda$, $\lambda'$, $\lambda''$ be the inclinations to the axis
of $x$ of the perpendiculars on the sides which touch these circles respectively; then
the equations of the three sides respectively are
\begin{align*}
(x-a)\cos\lambda + (y-b)\sin\lambda - c &= 0,\\
(x-a')\cos\lambda' + (y-b')\sin\lambda' - c' &= 0,\\
(x-a'')\cos\lambda'' + (y-b'')\sin\lambda'' - c'' &= 0.
\end{align*}
If the triangle be similar to a given triangle, then the differences of the angles
$\lambda$, $\lambda'$, $\lambda''$ will be given angles, or, what is the same thing, we may write
\[
\lambda = \theta + \xi, \quad \lambda' = \theta' + \xi, \quad \lambda'' = \theta'' + \xi,
\]
where $\theta$, $\theta'$, $\theta''$ are given angles, and $\xi$ is a variable angle. Let $\Delta$ be the area of
the triangle, then (disregarding a merely numerical factor) we have
\begin{multline*}
\sqrt{}\Delta = \sin(\lambda'-\lambda'')\,(a\cos\lambda + b\sin\lambda + c)\\
+ \sin(\lambda''-\lambda)\,(a'\cos\lambda' + b'\sin\lambda' + c') + \sin(\lambda-\lambda')\,(a''\cos\lambda'' + b''\sin\lambda'' + c'');
\end{multline*}
or, what is the same thing,
\begin{multline*}
\sqrt{}\Delta = \sin(\theta'-\theta'')\,\{a\cos(\theta+\xi) + b\sin(\theta+\xi) + c\}\\
+ \sin(\theta''-\theta)\,\{a'\cos(\theta'+\xi) + b'\sin(\theta'+\xi) + c'\}\\
+ \sin(\theta-\theta')\,\{a''\cos(\theta''+\xi) + b''\sin(\theta''+\xi) + c''\}.
\end{multline*}

It is now clear that the right-hand side will be independent of $\xi$, if only
\begin{multline*}
\sin(\theta'-\theta'')(a\cos\theta + b\sin\theta) + \sin(\theta''-\theta)(a'\cos\theta' + b'\sin\theta')\\
+ \sin(\theta-\theta')(a''\cos\theta'' + b''\sin\theta'') = 0,
\end{multline*}
\begin{multline*}
\sin(\theta'-\theta'')(-a\sin\theta + b\cos\theta) + \sin(\theta''-\theta)(-a'\sin\theta' + b'\cos\theta')\\
+ \sin(\theta-\theta')(-a''\sin\theta'' + b''\cos\theta'') = 0;
\end{multline*}
equations which show that, given the form of the triangle and the centres of two of
the circles, the centre of the third circle (in the porismatic case) is a determinate
unique point: and the theorem is thus proved.

\subsection*{[Vol.~\textsc{i}. pp.\ 137---141.]}

\paragraph*{1273.} (By the Editor [W.~J.~Miller, B.A.].)---In a given triangle let three
triangles be inscribed, by joining the points of contact of the inscribed circle, the
points where the bisectors of the angles meet the sides, and the points where the per-
pendiculars meet the sides; then will the corresponding sides of these three triangles
pass through the same point; also the triangle formed by the three points of inter-
section will be a circumscribed co-polar to the original triangle, and the pole will be
on the straight line in which the sides of the given triangle meet the bisectors of its
exterior angles.

\begin{center}
\textit{1. Solution by Professor Cayley.}
\end{center}

The theorem is, in fact, included in the following more general

\begin{center}
\textsc{Theorem.}
\end{center}

Let the points $O$, $O'$, $O''$, \ldots lie on a conic circumscribed about a
triangle $ABC$; then first the polars of the points $O$, $O'$, $O''$, \ldots in regard to the
triangle (see Note at the end of the Solution) pass through a fixed point $\Omega$. And
secondly, if by means of the point $O$, joining it with the vertices $A$, $B$, $C$, and taking
the intersections of these lines with the sides $BC$, $CA$, $AB$, respectively, we form a
triangle inscribed in the triangle $ABC$; and the like for the points $O'$, $O''$, \ldots; the
corresponding sides of the inscribed triangles meet in three points forming a triangle
circumscribed about the original triangle $ABC$, and such that the lines joining the
corresponding vertices of the last-mentioned two triangles meet in the point $\Omega$.

But, in order to see that the proposed theorem 1273 is in fact included under
the foregoing more general one, it is necessary to state the following

\begin{center}
\textsc{Subsidiary Theorem.}
\end{center}

Consider a conic inscribed in the triangle $ABC$, and
passing through the points $I$, $J$.

Take $O$ the pole of the line $IJ$ in regard to the conic; $O'$ the point of inter-
section of the lines joining the vertices of the triangle with the points of contact on
the opposite sides respectively; $O''$ the point of intersection of the lines $Al$, $Bm$, $Cn$,
where $l$ is a point on $BC$ such that the lines $lA$, $lBC$, $lI$, $lJ$ form a harmonic pencil,
and the like for the points $m$ and $n$ respectively.


Then the points $O$, $O'$, $O''$ lie on a conic circumscribed about the triangle $ABC$.

In fact, if in the subsidiary theorem the inscribed conic be a circle, and the
points $I$, $J$ be the circular points at infinity, the point $O$ will be the centre of the
circle, that is, the point of intersection of the interior bisectors of the angles; $O'$ will
be the point of intersection of the lines to the points of contact of the inscribed
circle; and $O''$ the point of intersection of the perpendiculars on the sides of the
triangle; and, these three points being on a conic circumscribed about the triangle,
the general theorem will apply to the three points in question.

I first prove the subsidiary theorem. Taking $x=0$, $y=0$, $z=0$ for the equations
of the sides of the triangle and $(\alpha, \beta, \gamma)$, $(\alpha', \beta', \gamma')$ for the coordinates of the points
$I$, $J$ respectively; the equation of the inscribed conic is
\[
\begin{vmatrix}
\sqrt{x}, & \sqrt{y}, & \sqrt{z} \\
\sqrt{\alpha}, & \sqrt{\beta}, & \sqrt{\gamma} \\
\sqrt{\alpha}', & \sqrt{\beta}', & \sqrt{\gamma}'
\end{vmatrix} = 0,
\]
or say
\[
a\sqrt{x} + b\sqrt{y} + c\sqrt{z} = 0,
\]
where
\[
a = \sqrt{\beta}\gamma' - \sqrt{\beta}'\gamma = p-p_1,\quad
b = \sqrt{\gamma}\alpha' - \sqrt{\gamma}'\alpha = q-q_1,\quad
c = \sqrt{\alpha}\beta' - \sqrt{\alpha}'\beta = r-r_1, \text{ suppose.}
\]
The coordinates of the point of intersection of the lines from the vertices to the
points of contact on the opposite sides are
\[
x : y : z = \frac{1}{a^2} : \frac{1}{b^2} : \frac{1}{c^2},
\]
that is,
\[
= \frac{1}{(p-p_1)^2} : \frac{1}{(q-q_1)^2} : \frac{1}{(r-r_1)^2}.
\]
The equation of the line $IJ$ is
\[
(\beta\gamma' - \beta'\gamma)\,x + (\gamma\alpha' - \gamma'\alpha)\,y + (\alpha\beta' - \alpha'\beta)\,z = 0;
\]
or, what is the same thing,
\[
(p^2 - p_1{}^2)\,x + (q^2 - q_1{}^2)\,y + (r^2 - r_1{}^2)\,z = 0.
\]
Representing this for a moment by $\lambda x + \mu y + \nu z = 0$, the coordinates of the pole of
this line, in regard to the inscribed conic $a\sqrt{x} + b\sqrt{y} + c\sqrt{z} = 0$, are as
\[
c^2\mu + b^2\nu : a^2\nu + c^2\lambda : b^2\lambda + a^2\mu.
\]
Now
\begin{align*}
c^2\mu + b^2\nu &= (r-r_1)^2(q^2 - q_1{}^2) + (q-q_1)^2(r^2 - r_1{}^2),\\
&= (r-r_1)(q-q_1)\,[(r-r_1)(q+q_1) + (q-q_1)(r+r_1)],\\
&= 2\,(r-r_1)(q-q_1)\,(qr - q_1r_1),
\end{align*}
and we have the like values for $a^2\nu + c^2\lambda$ and $b^2\lambda + a^2\mu$ respectively; hence, omitting
the symmetrical factor, we have, for the coordinates of the point in question,
\[
x : y : z = \frac{1}{pp_1} : \frac{1}{qq_1} : \frac{1}{rr_1}.
\]
Taking the equation of the line $Al$ to be $Qy + Rz = 0$, those of the lines $lI$, $lJ$ will be
\[
x = \lambda\,(Qy + Rz), \quad x = \lambda'\,(Qy + Rz),
\]
where
\[
\lambda = \frac{\alpha}{Q\beta + R\gamma}, \quad \lambda' = \frac{\alpha'}{Q\beta' + R\gamma'};
\]
and the harmonic condition gives $\lambda + \lambda' = 0$, that is,
\[
Q\,(\alpha\beta' + \alpha'\beta) + R\,(\alpha\gamma' + \alpha'\gamma) = 0;
\]
the equation of the line $Al$ is thus found to be
\[
(\gamma\alpha' + \gamma'\alpha)\,y = (\alpha\beta + \alpha'\beta)\,z;
\]
and, since we have the like forms for the equations of the lines $Bm$ and $Cn$, we have
for the coordinates of the point of intersection of these three lines
\[
x : y : z = \frac{1}{\beta\gamma' + \beta'\gamma} : \frac{1}{\gamma\alpha' + \gamma'\alpha} : \frac{1}{\alpha\beta' + \alpha'\beta},
\]
that is
\[
= \frac{1}{p^2 + p_1{}^2} : \frac{1}{q^2 + q_1{}^2} : \frac{1}{r^2 + r_1{}^2}.
\]
The equation of a conic circumscribed about the triangle $ABC$ is
\[
\frac{\lambda}{x} + \frac{\mu}{y} + \frac{\nu}{z} = 0,
\]
where $\lambda$, $\mu$, $\nu$ are arbitrary coefficients; and the condition for the three points being
in the conic is thus found to be
\[
\begin{vmatrix}
(p-p_1)^2, & (q-q_1)^2, & (r-r_1)^2 \\
pp_1, & qq_1, & rr_1 \\
p^2+p_1{}^2, & q^2+q_1{}^2, & r^2+r_1{}^2
\end{vmatrix} = 0,
\]
but, in virtue of the relations
\[
(p-p_1)^2 = -2pp_1 + (p^2 + p_1{}^2), \text{ \&c.,}
\]
this equation is identically true, and the subsidiary theorem is thus proved.

Passing now to the general theorem, I prove the first part of it as follows:

The equation of a conic circumscribed about the triangle $x=0$, $y=0$, $z=0$ is
\[
\frac{A}{x} + \frac{B}{y} + \frac{C}{z} = 0;
\]
hence, if $(\alpha, \beta, \gamma)$, $(\alpha', \beta', \gamma')$, $(\alpha'', \beta'', \gamma'')$ are the coordinates of any three points on
the conic, we have
\[
\frac{A}{\alpha} + \frac{B}{\beta} + \frac{C}{\gamma} = 0, \quad
\frac{A}{\alpha'} + \frac{B}{\beta'} + \frac{C}{\gamma'} = 0, \quad
\frac{A}{\alpha''} + \frac{B}{\beta''} + \frac{C}{\gamma''} = 0,
\]
and thence
\[
\begin{vmatrix}
\dfrac{1}{\alpha}, & \dfrac{1}{\beta}, & \dfrac{1}{\gamma} \\[1ex]
\dfrac{1}{\alpha'}, & \dfrac{1}{\beta'}, & \dfrac{1}{\gamma'} \\[1ex]
\dfrac{1}{\alpha''}, & \dfrac{1}{\beta''}, & \dfrac{1}{\gamma''}
\end{vmatrix} = 0,
\]
which is the condition for the intersection in a point of the three lines
\begin{align*}
\frac{x}{\alpha} + \frac{y}{\beta} + \frac{z}{\gamma} &= 0,\\
\frac{x}{\alpha'} + \frac{y}{\beta'} + \frac{z}{\gamma'} &= 0,\\
\frac{x}{\alpha''} + \frac{y}{\beta''} + \frac{z}{\gamma''} &= 0;
\end{align*}
and the theorem in question is thus proved. I remark, in passing, that the theorem
might also be stated as follows:---The locus of a point $O$, such that its polar in
regard to the triangle $ABC$ passes through a fixed point $\Omega$, is a conic circumscribed
about the triangle.

To prove the second part of the theorem, take for the coordinates of the points
$O$, $O'$, $O''$ respectively $(\alpha, \beta, \gamma)$, $(\alpha', \beta', \gamma')$, $(\alpha'', \beta'', \gamma'')$; then
\[
\begin{vmatrix}
\dfrac{1}{\alpha}, & \dfrac{1}{\beta}, & \dfrac{1}{\gamma} \\[1ex]
\dfrac{1}{\alpha'}, & \dfrac{1}{\beta'}, & \dfrac{1}{\gamma'} \\[1ex]
\dfrac{1}{\alpha''}, & \dfrac{1}{\beta''}, & \dfrac{1}{\gamma''}
\end{vmatrix} = 0,
\]
and if $X$, $Y$, $Z$ are the coordinates of the point $\Omega$, then we have
\begin{align*}
\frac{X}{\alpha} + \frac{Y}{\beta} + \frac{Z}{\gamma} &= 0,\\
\frac{X}{\alpha'} + \frac{Y}{\beta'} + \frac{Z}{\gamma'} &= 0,\\
\frac{X}{\alpha''} + \frac{Y}{\beta''} + \frac{Z}{\gamma''} &= 0.
\end{align*}
The equations of the sides of the inscribed triangle obtained by means of the
point $O$ are
\[
-\frac{x}{\alpha} + \frac{y}{\beta} + \frac{z}{\gamma} = 0, \qquad
\frac{x}{\alpha} - \frac{y}{\beta} + \frac{z}{\gamma} = 0, \qquad
\frac{x}{\alpha} + \frac{y}{\beta} - \frac{z}{\gamma} = 0,
\]
and the like for the triangles obtained by means of the points $O'$ and $O''$ respectively.
Hence, for a set of corresponding sides of the three triangles, we have, e.g.,
\[
-\frac{x}{\alpha} + \frac{y}{\beta} + \frac{z}{\gamma} = 0, \quad
-\frac{x}{\alpha'} + \frac{y}{\beta'} + \frac{z}{\gamma'} = 0, \quad
-\frac{x}{\alpha''} + \frac{y}{\beta''} + \frac{z}{\gamma''} = 0,
\]
and it is clear that these equations are simultaneously satisfied by the values
\[
x : y : z = -X : Y : Z,
\]
and we have the like expressions for the other sets of corresponding sides; that is,
we have for the coordinates of the vertices of the resulting triangle
\[
(-X : Y : Z), \quad (X : -Y : Z), \quad (X : Y : -Z);
\]
and hence also the equations of the sides of the triangle in question are
\[
\frac{y}{Y} + \frac{z}{Z} = 0, \quad \frac{z}{Z} + \frac{x}{X} = 0, \quad \frac{x}{X} + \frac{y}{Y} = 0,
\]
that is, it is a triangle circumscribed about the triangle $ABC$. The equations of the
lines joining the corresponding vertices of the two triangles are
\[
\frac{y}{Y} = \frac{z}{Z}, \quad \frac{z}{Z} = \frac{x}{X}, \quad \frac{x}{X} = \frac{y}{Y},
\]
and these lines meet in the point $(X : Y : Z)$, which is the point $\Omega$, the intersection
of the polars of $O$, $O'$, $O''$; the demonstration of the theorem is thus completed.

\{The expression Polar of a point in regard to a triangle denotes a line constructed
as follows:---viz., $O$ being the point and $ABC$ the triangle, then, taking on $BC$ a point
$a$, the harmonic in regard to the points $B$ and $C$ of the intersection of $BC$ by $AO$;
and in like manner on $CA$ and $AB$ the points $b$ and $c$ respectively, the three points
$a$, $b$, $c$ lie on a line which is the polar of the point $O$. If the equations of the
sides are $x=0$, $y=0$, $z=0$, and the coordinates of the point are $(\alpha, \beta, \gamma)$, then
the equation of the polar is $\dfrac{x}{\alpha} + \dfrac{y}{\beta} + \dfrac{z}{\gamma} = 0$; the equation may also be written
$(\alpha\delta_x + \beta\delta_y + \gamma\delta_z)^2\,xyz = 0$, and it thus appears that the line just defined as the polar is
in fact the second or line polar of the point in regard to the three lines $BC$, $CA$, $AB$
considered as forming a cubic curve.\}


\subsection*{[Vol.~\textsc{ii}. July to December 1864, pp.\ 6---9.]}

\paragraph*{1505.} (Proposed by Professor Cayley.)---If $P$, $Q$, $1$, $2$, $3$, $4$ be points on a conic,
then the four points $P1$, $Q2$; $P2$, $Q1$; $P3$, $Q4$; $P4$, $Q3$ lie on a conic passing through
the points $P$ and $Q$.

\begin{center}
\textit{Solution by the Proposer.}
\end{center}

This is an immediate consequence of the theorem of the anharmonic property of
the points of a conic. For if $(P1, P2, P3, P4)$ denote the anharmonic ratio of the
lines $P1$, $P2$, $P3$, $P4$, and so in other cases; then
\[
(P2, P1, P4, P3) = (P1, P2, P3, P4) = (Q1, Q2, Q3, Q4);
\]
that is
\[
(P2, P1, P4, P3) = (Q1, Q2, Q3, Q4),
\]
which proves the theorem.

In particular, if $P$, $Q$ are the circular points at infinity, then the conic is a circle.
Moreover the points $P1$, $Q2$; $P2$, $Q1$ are the antifocal points of $1$, $2$; viz., calling these
$1'$, $2'$, then $12$ and $1'2'$ are lines at right angles to each other, having a common
centre $O$, but such that $1'2' = i.12$, ($i=\surd-1$, as usual); or, what is the same thing,
$O1 = O2 = i.O1' = i.O2'$. And the theorem is as follows: viz., if $1$, $2$, $3$, $4$ are points
on a circle, and
\begin{align*}
&1',\ 2' \text{ are the antifocal points of } 1,\ 2,\\
&3',\ 4' \quad \text{``} \qquad\qquad \text{``} \qquad 3,\ 4,
\end{align*}
then $1'$, $2'$, $3'$, $4'$ are points on a circle.

As an \textit{\`{a} posteriori} proof, take the centre of the given circle as origin, so that
$(\alpha_1, \beta_1)$, $(\alpha_2, \beta_2)$, $(\alpha_3, \beta_3)$, $(\alpha_4, \beta_4)$ being the coordinates of $1$, $2$, $3$, $4$, and the radius
being taken as unity, we have
\[
\alpha_1{}^2 + \beta_1{}^2 = \alpha_2{}^2 + \beta_2{}^2 = \alpha_3{}^2 + \beta_3{}^2 = \alpha_4{}^2 + \beta_4{}^2 = 1.
\]
Suppose for a moment that $x$, $y$ are the coordinates of the antifocal points of $1$, $2$; we
have
\[
x - \alpha_1 \pm i\,(y - \beta_1) = 0, \quad x - \alpha_2 \mp i\,(y - \beta_2) = 0,
\]
that is
\[
x + iy = \alpha_1 + i\beta_1, \quad x - iy = \alpha_2 - i\beta_2,
\]
for the coordinates of the one point; and similarly
\[
x - iy = \alpha_1 - i\beta_1, \quad x + iy = \alpha_2 + i\beta_2,
\]
for the coordinates of the other point.

Hence, taking the new coordinates
\[
X = x + iy, \quad Y = x - iy,
\]
and similarly $A_1 = \alpha_1 + i\beta_1$, $B_1 = \alpha_1 - i\beta_1$, \&c.; the coordinates of the antifocal points $1'$, $2'$
are $(A_1, B_2)$ and $(A_2, B_1)$ respectively; but we have $A_1B_1 = \alpha_1{}^2 + \beta_1{}^2 = 1$, $A_2B_2 = \alpha_2{}^2 + \beta_2{}^2 = 1$;
so that $B_1 = \dfrac{1}{A_1}$, $B_2 = \dfrac{1}{A_2}$; and the coordinates are $\left(A_1, \dfrac{1}{A_2}\right)$, $\left(A_2, \dfrac{1}{A_1}\right)$ respectively.
Similarly the coordinates of the antifocal points $(3', 4')$ are $\left(A_3, \dfrac{1}{A_4}\right)$, $\left(A_4, \dfrac{1}{A_3}\right)$ respec-
\noindent tively.

Take as the equation of the circle through the two pairs of antifocal points
\[
x^2 + y^2 + 2\lambda x + 2\mu y + \nu = 0,
\]
or, what is the same thing,
\[
XY + \lambda\,(X+Y) - i\mu\,(X-Y) + \nu = 0,
\]
that is
\[
XY + LY + MX + N = 0,
\]
if
\[
L = \lambda + i\mu, \quad M = \lambda - i\mu, \quad N = \nu.
\]
We ought then to have
\begin{align*}
\frac{A_1}{A_2} + L\,\frac{1}{A_2} + MA_1 + N &= 0,\\
\frac{A_2}{A_1} + L\,\frac{1}{A_1} + MA_2 + N &= 0,\\
\frac{A_3}{A_4} + L\,\frac{1}{A_4} + MA_3 + N &= 0,\\
\frac{A_4}{A_3} + L\,\frac{1}{A_3} + MA_4 + N &= 0;
\end{align*}
and these will exist simultaneously, if
\[
\begin{vmatrix}
\dfrac{A_1}{A_2}, & \dfrac{1}{A_2}, & A_1, & 1 \\[1ex]
\dfrac{A_2}{A_1}, & \dfrac{1}{A_1}, & A_2, & 1 \\[1ex]
\dfrac{A_3}{A_4}, & \dfrac{1}{A_4}, & A_3, & 1 \\[1ex]
\dfrac{A_4}{A_3}, & \dfrac{1}{A_3}, & A_4, & 1
\end{vmatrix} = 0,
\]
an identical equation which is easily verified. It, in fact, gives
\begin{multline*}
\left(\frac{1}{A_2} - \frac{1}{A_1}\right)(A_3 - A_4) - (A_1 - A_2)\left(\frac{1}{A_4} - \frac{1}{A_3}\right) + \left(\frac{A_1}{A_2} - \frac{A_2}{A_1}\right)(1-1) +\\
(1-1)\left(\frac{A_3}{A_4} - \frac{A_4}{A_3}\right) - \left(\frac{1}{A_3} - \frac{1}{A_1}\right)(A_3 - A_4) + (A_1 - A_2)\left(\frac{1}{A_4} - \frac{1}{A_3}\right) = 0,
\end{multline*}
which is obviously true. The equation may also be written
\[
\begin{vmatrix}
1, & A_1, & A_2, & A_1A_2 \\
1, & A_2, & A_1, & A_1A_2 \\
1, & A_3, & A_4, & A_3A_4 \\
1, & A_4, & A_3, & A_3A_4
\end{vmatrix} = 0,
\]
and in this form it expresses the known theorem of the equality of the anharmonic
ratios of $(A_1, A_2, A_3, A_4)$ and $(A_2, A_1, A_4, A_3)$.

But, in order to actually find the circle, we may write
\begin{align*}
XY + LY + MX \phantom{+N\;} &= 0,\\
A_1 + L \phantom{+MA_1} + MA_1A_2 + NA_2 &= 0,\\
A_2 + L \phantom{+MA_2} + MA_1A_2 + NA_1 &= 0,\\
A_3 + L \phantom{+MA_3} + MA_3A_4 + NA_4 &= 0,
\end{align*}
and eliminating $L$, $M$, $N$, the equation of the circle is
\[
\begin{vmatrix}
XY, & Y, & X, & 1 \\
A_1, & 1, & A_1A_2, & A_2 \\
A_2, & 1, & A_1A_2, & A_1 \\
A_3, & 1, & A_3A_4, & A_4
\end{vmatrix} = 0,
\]
or, reducing, this is
\begin{multline*}
(A_2-A_1)\,[\,XY\,(A_3A_4 - A_1A_2) + Y\,\{A_1A_2\,(A_3+A_4) - A_3A_4\,(A_1+A_2)\}\\
+ X\,(A_1+A_2-A_3-A_4) + (A_3A_4 - A_1A_2)\,] = 0,
\end{multline*}
or say
\begin{multline*}
XY\,(A_1A_2 - A_3A_4) + Y\,\{A_3A_4\,(A_1+A_2) - A_1A_2\,(A_3+A_4)\}\\
+ X\,\{A_3+A_4-A_1-A_2\} + (A_1A_2 - A_3A_4) = 0;
\end{multline*}
that is
\[
\begin{vmatrix}
XY+1, & X, & Y \\
A_1+A_2, & A_1A_2, & 1 \\
A_3+A_4, & A_3A_4, & 1
\end{vmatrix} = 0,
\]
which is the required equation; or, transforming to the original axes, we have $x+iy=X$,
$x-iy=Y$, \&c., and therefore $XY=x^2+y^2$; and the equation becomes
\[
\begin{vmatrix}
x^2+y^2+1, & x+iy, & x-iy \\
\alpha_1+\alpha_2+i(\beta_1+\beta_2), & (\alpha_1+i\beta_1)(\alpha_2+i\beta_2), & 1 \\
\alpha_3+\alpha_4+i(\beta_3+\beta_4), & (\alpha_3+i\beta_3)(\alpha_4+i\beta_4), & 1
\end{vmatrix} = 0,
\]
which is the equation of the circle through the two pairs of antifocal points.

\{Note. The \textit{second} form of the equation of the circle may be otherwise deduced
from the \textit{first}, without expanding the determinants, by the following method:
\[
\begin{vmatrix}
XY, & Y, & X, & 1 \\
A_1, & 1, & A_1A_2, & A_2 \\
A_2, & 1, & A_1A_2, & A_1 \\
A_3, & 1, & A_3A_4, & A_4
\end{vmatrix}
=
\begin{vmatrix}
XY+1, & Y, & X, & 1 \\
A_1+A_2, & 1, & A_1A_2, & A_2 \\
A_1+A_2, & 1, & A_1A_2, & A_1 \\
A_3+A_4, & 1, & A_3A_4, & A_4
\end{vmatrix}
=
\]
\[
\begin{vmatrix}
XY+1, & Y, & X, & 1 \\
A_1+A_2, & 1, & A_1A_2, & A_2 \\
0, & 0, & 0, & A_1-A_2 \\
A_3+A_4, & 1, & A_3A_4, & A_4
\end{vmatrix}
= (A_1-A_2)
\begin{vmatrix}
XY+1, & X, & Y \\
A_1+A_2, & A_1A_2, & 1 \\
A_3+A_4, & A_3A_4, & 1
\end{vmatrix};
\]
therefore
\[
\begin{vmatrix}
XY+1, & X, & Y \\
A_1+A_2, & A_1A_2, & 1 \\
A_3+A_4, & A_3A_4, & 1
\end{vmatrix} = 0.\quad\text{Ed. [W. J. M.]}
\]

\}


\subsection*{[Vol.~\textsc{ii}. pp.\ 22---24.]}

\paragraph*{1513.} (Proposed by the Rev.~J.~Blissard, B.A.)---Prove the following formul\ae:
\begin{align*}
&(1)\quad \frac{(x-1)(x-2)\ldots(x-n)}{x(x+1)\ldots(x+n-1)} = \\
&1 + (-1)^n\left\{n.\frac{1}{x} - \frac{n(n^2-1^2)}{1^2}.\frac{1}{x+1} + \frac{n(n^2-1^2)(n^2-2^2)}{1^2.2^2}.\frac{1}{x+2} - \&\text{c.}\right\}.
\end{align*}

(2) The above formula expressed as
\[
\frac{(\Gamma x)^2}{\Gamma(x-n)\,\Gamma(x+n)} = 1 - \frac{n^2}{1}.\frac{1}{x} + \frac{n^2(n^2-1^2)}{1.2}.\frac{1}{x(x+1)} - \frac{n^2(n^2-1^2)(n^2-2^2)}{1.2.3}.\frac{1}{x(x+1)(x+2)} + \&\text{c.}
\]
and show that this equation is subject to the sole restriction that when $n$ is not
integral $x$ must not be negative.

\begin{center}
\textit{Solution by Professor Cayley; and X.~U.~J.}
\end{center}

Let $n$ be a positive integer, and suppose that $[x]^n$ denotes as usual the factorial
$x(x-1)\ldots(x-n+1)$; then we have
\begin{align*}
[x+k]^n &= (1+\Delta)^k\,[x]^n = \left(1 + k\Delta + \frac{k(k-1)}{1.2}\Delta^2 + \&\text{c.}\right)[x]^n\\
&= [x]^n + \frac{kn}{1}\,[x]^{n-1} + \frac{k(k-1)\,n(n-1)}{1.2}\,[x]^{n-2} + \&\text{c.};
\end{align*}
or putting $k=-n$ we have
\[
[x-n]^n = [x]^n - \frac{n^2}{1}\,[x]^{n-1} + \frac{n^2(n^2-1^2)}{1.2}\,[x]^{n-2} - \&\text{c.}
\]
Writing herein $(x+n-1)$ for $x$, and dividing by $[x+n-1]^n$, we have
\[
\frac{[x-1]^n}{[x+n-1]^n} = 1 - \frac{n^2}{1}.\frac{1}{x} + \frac{n^2(n^2-1^2)}{1.2}.\frac{1}{x(x+1)} - \&\text{c.};
\]
or, what is the same thing,
\[
\frac{(\Gamma x)^2}{\Gamma(x-n)\,\Gamma(x+n)} = 1 - \frac{n^2}{1}.\frac{1}{x} + \frac{n^2(n^2-1^2)}{1.2}.\frac{1}{x(x+1)} - \&\text{c.},
\]
which is the formula (2). The foregoing demonstration applies to the case of $n$ a
positive integer; but as the two sides are respectively unaltered when $n$ is changed
into $-n$, it is clear that the formula holds good also for $n$ a negative integer. The
right hand side is the hypergeometric series $F(n, -n, x, 1)$ and the formula therefore is
\[
\frac{(\Gamma x)^2}{\Gamma(x-n)\,\Gamma(x+n)} = F(n, -n, x, 1),
\]
a particular case of the known formula
\[
\frac{\Gamma(\gamma)\,\Gamma(\gamma-\alpha-\beta)}{\Gamma(\gamma-\alpha)\,\Gamma(\gamma-\beta)} = F(\alpha, \beta, \gamma, 1),
\]
which when $\alpha$ or $\beta$ is a positive integer is a mere identity, true therefore for all
values of $\gamma$; but if neither $\alpha$ nor $\beta$ is a positive integer, then the right hand side
is an infinite series which is only convergent for $\gamma > \alpha + \beta$. In the particular case we
have $\alpha=n$, $\beta=-n$, $\gamma=x$; hence if $n$ be a positive or negative integer, the formula
is an identity, but if $n$ be fractional, the condition of convergence is $x > 0$, that is,
$x$ must be positive.

To prove the formula (1) it is only necessary to remark, that ($n$ being a positive
integer) the quantity $\dfrac{[x-1]^n}{[x+n-1]^n}$ is a rational fraction, the numerator and denominator
whereof are of the same degree $n$, and which becomes $=1$ for $x=\infty$. Hence, decom-
\noindent posing it into simple fractions, we may write
\[
\frac{[x-1]^n}{[x+n-1]^n} = 1 + S_r.\frac{A_r}{x+r},
\]
where the summation extends from $r=0$ to $r=n-1$ both inclusive. And we have
\[
A_r = \left\{\frac{(x+r)\,[x-1]^n}{[x+n-1]^n}\right\}_{x=-r},
\]
or, observing that $[x+n-1]^n = [x+n-1]^{n-r-1}\,(x+r)\,[x+r-1]^r$, we have
\begin{align*}
A_r &= \left\{\frac{[x-1]^n}{[x+n-1]^{n-r-1}\,[x+r-1]^r}\right\}_{x=-r} = \frac{[-r-1]^n}{[n-r-1]^{n-r-1}\,[-1]^r}\\
&= \frac{(-)^n\,[n+r]^n}{[n-r-1]^{n-r-1}\,(-)^r\,[r]^r} = (-)^{n+r}\,\frac{[n+r]^{n+r}}{[n-r-1]^{n-r-1}\,[r]^r\,[r]^r} = (-)^{n+r}\,\frac{[n+r]^{2r+1}}{[r]^r\cdot[r]^r}.
\end{align*}
Hence the formula is
\[
\frac{[x-1]^n}{[x+n-1]^n} = 1 + (-)^n\,S_r\,(-)^r\,\frac{[n+r]^{2r+1}}{[r]^r\,[r]^r}\,\frac{1}{x+r},
\]
or, as this may also be written,
\[
\frac{(x-1)(x-2)\ldots(x-n)}{x(x+1)\ldots(x+n-1)} = 1 + (-)^n\left\{n.\frac{1}{x} - \frac{n(n^2-1^2)}{1^2}.\frac{1}{x+1} + \frac{n(n^2-1^2)(n^2-2^2)}{1^2.2^2}.\frac{1}{x+2} - \&\text{c.}\right\},
\]
which is the formula in question.

\subsection*{[Vol.~\textsc{ii}. pp.\ 51, 52.]}

\paragraph*{1512.} (Proposed by Professor Cayley.)---It is possible to construct a hexagon
$123456$, inscribed in a conic, and such that the diagonals $14$, $25$, $36$ pass respectively
through the Pascalian points (intersections of opposite sides) $23$, $56$; $34$, $61$; $45$, $12$.
Given the points $1$, $2$; $4$, $5$; to construct the hexagon.

\begin{center}
\textit{Solution by the Proposer.}
\end{center}

Let $12$, $45$, meet in $O$, and through $O$ draw at pleasure a line meeting $14$ in $P$,
and $25$ in $Q$; let $P2$, $Q4$ meet in $3$, and $P5$, $Q1$ in $6$; then the line $36$ will pass
through $O$, and this being so, the hexagon $123456$ satisfies the required conditions.

We have to show that $36$ passes through $O$. Let $Q4$ meet $O12$ in $A$, and $P2$
meet $O45$ in $B$; then the points $6$, $3$, $O$, are the intersections of corresponding sides
of the triangles $A1Q$, $B5P$; and in order that these points may lie in a line, the
lines joining the corresponding vertices must meet in a point, that is, we have to show
that the lines $15$, $AB$, $PQ$ meet in a point. The property is in fact as follows; viz.,
given the points $2$, $4$; and also the points $Q$, $O$, $P$ lying in a line; then constructing
the points $1$, $5$, $A$, $B$, which are the respective intersections of $P4$, $O2$; $Q2$, $O4$;
$Q4$, $O2$; $P2$, $O4$; the lines $15$, $AB$, $PQ$ will meet in a point. Take $x=0$, $y=0$,
$z=0$ for the respective equations of $P2$, $Q4$, $PQ$; then $O$ is an arbitrary point in the
line $PQ$, say that for the point $O$ we have $z=0$, $ax+by=0$; also $O2$, $O4$ are
arbitrary lines through $O$: say that their equations are $ax+by+\lambda z=0$; $ax+by+\mu z=0$;
then we have for the points $A$ and $B$, respectively, $ax+by+\mu z=0$, $y=0$; $ax+by+\mu z=0$,
$x=0$; hence the equation of $AB$ is $\mu ax + \lambda by + \lambda\mu z = 0$. The equation of $P4$ is
$ax+\mu z=0$, and that of $Q2$ is $by+\lambda z=0$; the point $1$ is therefore given by
$ax+\mu z=0$, $ax+by+\lambda z=0$; and $5$ by $by+\lambda z=0$, $ax+by+\mu z=0$; hence the equation
of $15$ is $\mu ax + \lambda by + (\mu^2 - \mu\lambda + \lambda^2)\,z = 0$; and the equation of $PQ$ being $z=0$, it is
clear that the three lines $AB$, $15$, $PQ$ intersect in the point given by the equations
$\mu ax + \lambda by = 0$, $z=0$.

\textsc{Obs.~1.} By inspection of the figure we see that $3PQ$ is a triangle whereof the
sides $3P$, $3Q$, $PQ$ pass respectively through the fixed points $2$, $4$, $O$; while the vertices
$P$ and $Q$ lie in the fixed lines $14$, $25$; the locus of the vertex $3$ is consequently a
conic; and the like as regards the triangle $6PQ$.

\textsc{Obs.~2.} The regular hexagon projects into a hexagon inscribed in a conic and
circumscribed about another conic having double contact therewith; in the hexagon
so obtained (as appears at once by the consideration of the regular hexagon) the
above-mentioned property holds; but the in-and-circumscribed hexagon has the additional
property that the three diagonals meet in a point, and it is therefore a less general
figure than the hexagon of the foregoing theorem. It would, I think, be worth while
to study further the hexagon of the theorem.

\{Note. In the solution of Question 1548 it is shown that if two pairs of
opposite sides of \textit{any hexagon} intersect each on a diagonal produced, so likewise will
the third pair.

A slight variation of Professor Cayley's proof may be obtained by finding the
equations of $P5$, $Q1$, and thence of $36$, which are respectively
\[
ax - (\lambda - \mu)\,z = 0, \quad bx + (\lambda - \mu)\,z = 0, \quad ax + by = 0,
\]
showing that $36$ passes through $O$. Ed. [W.~J.~M.].\}

\subsection*{[Vol.~\textsc{ii}. pp.\ 70---72.]}

\paragraph*{1562.} (Proposed by F.~D.~Thomson, M.A.)---Find the locus of the points of contact
of tangents drawn from a given point to a conic circumscribing a given quadrangle.
The quadrangle being supposed convex, trace the changes of form of the locus for
different positions of the given point.

\begin{center}
\textit{Solution by Professor Cayley; and the Proposer.}
\end{center}

Let $O$ be the given point; $1$, $2$, $3$, $4$ the vertices of the given quadrangle;
$A$, $B$, $C$ the centres of the quadrangle, viz., $A$ the intersection of the lines $14$, $23$;
$B$ of $24$, $31$; $C$ of $34$, $12$. The polars of $O$ in regard to the several circumscribed
conics intersect in a point $O'$. This being so, the locus is a cubic passing through
the nine points $1$, $2$, $3$, $4$, $A$, $B$, $C$, $O$, $O'$, and which is moreover such that the
tangents at the four points $1$, $2$, $3$, $4$ meet the cubic in the point $O$, and the
tangents at the four points $A$, $B$, $C$, $O$ meet the cubic in the point $O'$. It is to be
remarked that the nine points are so related to each other that a cubic through
any eight of these points passes through the remaining ninth point; say a cubic
through $1$, $2$, $3$, $4$, $A$, $B$, $C$, $O$ passes through $O'$; the nine points consequently do
not determine the cubic; but the cubic will be determined, e.g., by the conditions
that it passes through $1$, $2$, $3$, $4$, $A$, $B$, $C$, $O$, and has $O1$ for the tangent at $1$.
The series of cubics corresponding to different positions of the point $O$ is identical
with the series of cubics passing through the seven points $1$, $2$, $3$, $4$, $A$, $B$, $C$
(which is the series of cubics through the centres of a quadrangle).

Conversely any given cubic curve may be taken to be a cubic of the series; and
the points $1$, $2$, $3$, $4$ will then be determined as follows, viz., $1$, $2$, $3$, $4$ are the points
of contact of the tangents to the cubic from an arbitrary point $O$ on the cubic; and
then taking as before $A$, $B$, $C$ for the intersections of $14$, $23$, of $24$, $31$ and of
$34$, $12$, respectively, the points $A$, $B$, $C$ will lie on the cubic, and the tangents at
$A$, $B$, $C$, $O$ will meet the cubic in a point $O'$. I call to mind that a cubic curve
without singularities is either \textit{complex} or \textit{simplex}; in the simplex kind there can be
drawn from any point of the curve two, and only two, real tangents to the curve;
in the complex kind, there can be drawn four real tangents or else no real tangent,
viz.\ from any point on a certain branch of the curve there can be drawn four real
tangents, from a point on the remaining portion of the curve no real tangent.
Hence, in the foregoing construction, in order that the points $1$, $2$, $3$, $4$ may be real,
the given cubic must be of the complex kind, and the point $O$ must be taken on
the branch which has through each of its points four real tangents.

The foregoing results may be established \textit{geometrically} or \textit{analytically}; but for
brevity I merely indicate the analytical demonstration. Suppose first, that the points
$1$, $2$, $3$, $4$ are given as the intersections of the conics $U=0$, $V=0$; let $\alpha$, $\beta$, $\gamma$ be
the coordinates of the point $O$, and write $D=\alpha\delta_x + \beta\delta_y + \gamma\delta_z$, so that $DU=0$ and
$DV=0$ are the equations of the polars of $O$ in regard to the conics $U=0$, $V=0$
respectively. The equation of any conic through the four points is $U+kV=0$; and
the equation of the polar of $O$ in regard thereto is $DU+kDV=0$; eliminating $k$
from these equations, we have $UDV - VDU = 0$, which is the equation of the given
locus. We see at once that it is a cubic curve passing through the points
$(U=0,\ V=0)$, that is, the points $1$, $2$, $3$, $4$; and through the point $DU=0$, $DV=0$,
that is, the point $O'$; it also follows without difficulty that the curve passes through
the point $O$. But for the remaining results it is better to particularize the conics
$U=0$, $V=0$. Let the equations of $12$, $23$, $34$, $41$ be $x=0$, $y=0$, $z=0$, $w=0$
respectively, (where $x+y+z+w=0$); and in the same system, let $\alpha$, $\beta$, $\gamma$, $\delta$ be the
coordinates of $O$ ($\alpha+\beta+\gamma+\delta=0$), then $xz=0$, $yw=0$ are each of them a conic
(pair of lines) passing through the four points; and we may therefore write $U=yw$,
$V=xz$; the equation $UDV-VDU=0$ thus becomes $yw\,(\alpha z + \gamma x) - xz\,(\beta w + \delta y) = 0$, or,
as this equation may also be written,
\[
\frac{\alpha}{x} - \frac{\beta}{y} + \frac{\gamma}{z} - \frac{\delta}{w} = 0,
\]
which is the equation of the cubic curve; and from this form the several above-
mentioned results may be obtained without difficulty.

To give an idea of the form of the curve corresponding to a given convex
quadrangle $1234$, and given position of the point $O$, I suppose that $O$ is situate
\textit{within} the quadrangle, for instance in the triangle $B12$. The mere inspection of the
figure, and consideration of the conditions which are to be satisfied by the cubic curve,
is enough to show that this is of the form described by Newton as \textit{angui\kern0pt nea cum
ovali}, viz., the oval passes through the points $3$, $4$, $A$, $B$, and the serpentine branch
through the points $1$, $2$, $C$, $O$, $O'$. But the complete discussion of the different cases
would be somewhat laborious.

[A geometrical investigation of the locus is given on p.~124 of Cremona's \textit{Teoria
Geometrica delle Curve Piane}. Ed. [W.~J.~M.]]


\subsection*{[Vol.~\textsc{ii}. pp.\ 89, 90.]}

\paragraph*{1533.} (Proposed by Professor Cayley.)---If on the sides of a triangle there are
taken three points, one on each side; and if through the three points and the three
vertices of the triangle there are drawn a cubic curve and a quartic curve, inter-
secting in six other points; then there exists a quintic curve passing through each
of the three points, and having each of the six points for a double point.

\begin{center}
\textit{Solution by the Proposer.}
\end{center}

Let $P=0$ be the equation of the quartic curve, $Q=0$ the equation of the cubic
curve, $M=0$ the equation of the three sides of the triangle; then if we can find
$A$, $B$, $C$ functions of the orders $0$, $1$, $2$ respectively, and $U$ a function of the fifth
order, such that we have identically $MU = AP^2 + BPQ + CQ^2$; we have $MU=0$, a curve
of the eighth order, having a double point at each of the points $(P=0,\ Q=0)$, which
points are the three vertices of the triangle, the three points, and the six points;
but the curve $MU=0$ is made up of the curve $M=0$ (the three sides of the
triangle, being a cubic curve having each of the vertices for a double point, and
passing through each of the three points) and of a certain quintic curve $U=0$;
hence the quintic curve must pass through each of the three points, and have a
double point at each of the six points; or there exists a quintic curve satisfying the
conditions of the theorem.

I take $x=0$, $y=0$, $z=0$ for the equations of the three sides of the triangle, and
then (the constants being all of them arbitrary) writing for shortness
\begin{alignat*}{3}
\xi &= .\phantom{''}by+cz, \qquad & X &= .\phantom{''}\beta y+\gamma z, \qquad & \Theta &= \lambda x+\mu y+\nu z,\\
\eta &= a'x.\phantom{+c'z} + c'z, \qquad & Y &= \alpha'x.\phantom{+\gamma'z} + \gamma'z, \qquad & \\[-0.5ex]
\zeta &= a''x+b''y.\phantom{+c''z}, \qquad & Z &= \alpha''x+\beta''y.\phantom{+\gamma''z}, \qquad &
\end{alignat*}
I assume that the three points are given by the equations $(x=0,\ \xi=0)$, $(y=0,\ \eta=0)$,
$(z=0,\ \zeta=0)$, respectively. This being so, we may write
\[
Q = yz\xi\delta + zx\eta\delta' + xy\zeta\delta'' + xyz\epsilon = 0, \quad
-P = yz\xi X + zx\eta Y + xy\zeta Z + xyz\Theta = 0,
\]
for the equations of the cubic curve and the quartic curve respectively. We have
of course $M=xyz=0$ for the equation of the three sides of the triangle, and the
identity to be satisfied is $xyzU = AP^2 + BPQ + CQ^2$.

I was led to the values of $A$, $B$, $C$ by considerations founded on the theory of
curves in space. We have
\begin{align*}
A &= \delta\delta'\delta'',\\
B &= (\delta'\alpha'' + \delta''\alpha)\,\delta x + (\delta''\beta + \delta\beta'')\,\delta' y + (\delta\gamma' + \delta'\gamma)\,\delta'' z,\\
C &= \alpha'\alpha''\delta x^2 + \beta''\beta\delta' y^2 + \gamma\gamma'\delta'' z^2 + (\gamma\beta'\delta' + \gamma'\beta\delta'')\,yz + (\alpha'\gamma\delta'' + \alpha''\gamma'\delta)\,zx + (\beta''\alpha\delta + \beta\alpha'\delta')\,xy;
\end{align*}
and with these values it is easy to show that the function $AP^2 + BPQ + CQ^2$ contains
the factor $xyz$; for substituting the values of $P$, $Q$, all the terms of $AP^2 + BPQ + CQ^2$
contain explicitly the factor $xyz$, except the terms
\begin{multline*}
A\,(y^2z^2\xi^2X^2 + z^2x^2\eta^2Y^2 + x^2y^2\zeta^2Z^2) - B\,(y^2z^2\xi^2X\delta + z^2x^2\eta^2Y\delta' + x^2y^2\zeta^2Z\delta'')\\
+ C\,(y^2z^2\xi^2\delta^2 + z^2x^2\eta^2\delta'^2 + x^2y^2\zeta^2\delta''^2);
\end{multline*}
and these terms will contain the factor $xyz$, if only the expressions $AX^2 - BX\delta + C\delta^2$,
$AY^2 - BY\delta' + C\delta'^2$, $AZ^2 - BZ\delta'' + C\delta''^2$ contain respectively the factors $x$, $y$, $z$. But
$AX^2 - BX\delta + C\delta^2$ will contain the factor $x$, if only the expression vanishes for $x=0$;
and for $x=0$ we have
\begin{multline*}
AX^2 - BX\delta + C\delta^2 = 0 =\\
\delta\delta'\delta''\,(\beta y + \gamma z)^2 - [\delta'\delta''\,(\beta y + \gamma z) + \delta\,(\beta''\delta' y + \gamma'\delta'' z)]\,\delta\,(\beta y + \gamma z) + (\beta y + \gamma z)\,(\beta''\delta' y + \gamma'\delta'' z)\,\delta^2;
\end{multline*}
that is, $AX^2 - BX\delta + C\delta^2$ contains the factor $x$; and by symmetry the other two
expressions contain the factors $y$ and $z$ respectively. The excepted terms contain therefore
the factor $xyz$; and there exists therefore a quintic function $U=(AP^2+BPQ+CQ^2)\div xyz$;
which proves the theorem.

The values of $A$, $B$, $C$ were obtained by considering the surface $w=\dfrac{P}{Q}$, which,
as is at once seen, contains upon itself the three lines
\[
\left(x=0,\ w=-\frac{X}{\delta}\right),\quad
\left(y=0,\ w=-\frac{Y}{\delta'}\right),\quad
\left(z=0,\ w=-\frac{Z}{\delta''}\right)
\]
or as these equations may be written
\begin{align*}
(x=0, \phantom{y=0,} &\beta y + \gamma z + \delta w = 0),\\
(y=0, \alpha'x \phantom{+\gamma'z} &+ \gamma' z + \delta' w = 0),\\
(z=0, \alpha''x + \beta''y \phantom{+\delta''w} &+ \delta'' w = 0);
\end{align*}
and then seeking for the equation of the hyperboloid which passes through the three
lines, this is found to be $Aw^2 + Bw + C = 0$, where $A$, $B$, $C$ have the before-mentioned
values.

If in the foregoing theorem the cubic is considered as a given cubic curve, and
the three points as three arbitrary points on the cubic, the question then arises to
find the triangle; or we have the problem proposed as Question 1607.

\subsection*{[Vol.~\textsc{ii}. p.\ 91.]}

\paragraph*{1542.} (Proposed by Professor Cayley.)---If a given line meet two given conics
in the points $(A, B)$ and $(A', B')$ respectively; and if $(A'', B'')$ be the sibi-conjugate
points (or foci) of the pairs $(A, A')$ and $(B, B')$, or of the pairs $(A, B')$ and $(A', B)$,
then $(A'', B'')$ lie on a conic passing through the four points of intersection of the
two given conics.

\subsection*{[Vol.~\textsc{ii}. pp.\ 97---100.]}

\paragraph*{1606.} (Proposed by the Editor, [W.~J.~M.].)---Solve the following problems:

$(\alpha)$ Through three given points to draw a conic whose foci shall lie in two given
lines.

$(\beta)$ Through four given points to draw a conic such that one of its chords of
intersection with a given conic shall pass through a given point.

$(\gamma)$ Through two given points to draw a circle such that its chords of inter-
section with a given circle shall pass through a given point.

\begin{center}
\textit{Solution by Professor Cayley.}
\end{center}

$(\alpha)$ Through three given points to draw a conic whose foci shall lie in two given
lines.

The focus of a conic is a point such that the lines joining it with the two
circular points at infinity (say the points $I$, $J$) are tangents to the conic. Hence the
question is, in a given line to find a point $A$, and in another given line to find a
point $B$, such that there exists a conic touching the four lines $AI$, $AJ$, $BI$, $BJ$
(where $I$, $J$ are any given points) and besides passing through three given points.
More generally, instead of the lines from $A$, $B$ to the given points $I$, $J$, we may
consider the tangents from $A$, $B$, to a given conic $\Theta$; the question then is, in a
given line to find a point $A$, and in another given line to find a point $B$, such that
there exists a conic touching the tangents from $A$, $B$ to a given conic $\Theta$, and besides
passing through three given points. It is rather more convenient to consider the
reciprocal question---though it is to be borne in mind that for any two reciprocal
questions a solution of the one question by means of coordinates $(x, y, z)$ regarded as
point-coordinates is in fact a solution of the other question by means of the same
coordinates $(x, y, z)$ regarded as line-coordinates. The reciprocal question is: through
a given point to draw a line $A$, and through another given point to draw a line $B$,
such that there exists a conic passing through the intersections of these lines with a
given conic $\Theta$, and besides touching three given lines. The given points may be taken
to be $(x=0,\ z=0)$, $(y=0,\ z=0)$; this determines the line $z=0$, but not the lines
$x=0$, $y=0$, so that the point $(x=0,\ y=0)$ may without loss of generality be supposed
to lie on the conic $\Theta$; the equation of this conic will therefore be
\[
(a, b, 0, f, g, h\!\!\between\! x, y, z)^2 = 0.
\]
I take $\alpha z + \gamma x = 0$ for the equation of the line $A$, $\mu y + \nu z = 0$ for the equation of the
line $B$ (so that the quantities to be determined are the ratios $\alpha:\gamma$ and $\mu:\nu$); this
being so, the required conic passes through the intersections of these lines with the
conic $\Theta$; its equation will therefore be
\[
(a, b, 0, f, g, h\!\!\between\! x, y, z)^2 + 2(\alpha z + \gamma x)(\mu y + \nu z) = 0;
\]
or what is the same thing
\[
(a, b, 2\nu\gamma, f+\mu\gamma, g+\nu\alpha, h+\mu\alpha\!\!\between\! x, y, z)^2 = 0;
\]
where $\alpha$, $\gamma$, $\mu$, $\nu$ have to be determined in such manner that this conic may touch
three given lines. It is to be observed that $\alpha$, $\gamma$, $\mu$, $\nu$ enter into the equation
through the combinations $\alpha\mu$, $\alpha:\gamma$, and $\mu:\nu$, so that there are really only three
disposable quantities.

\subsection*{[Vol.~\textsc{iii}. January to July, 1865, p.\ 29.]}

\paragraph*{1607.} (Proposed by Professor Cayley.)---In a given cubic curve to inscribe a
triangle such that the three sides shall pass respectively through three given points on
the curve.

\subsection*{[Vol.~\textsc{iii}. pp.\ 60---63.]}

\paragraph*{1647.} (Proposed by Professor Cayley.)---Find the locus of the foci of an ellipse
of given major axis, passing through three given points.

\{In connexion with the problem the Proposer remarks as follows:

Let $A$, $B$, $C$ be the given points; take $P$ an arbitrary point (not in general in
the plane of the three given points), then we may find a point $Q$ (not in general in
the plane of the three given points) such that $QA+AP = QB+BP = QC+CP =$
given major axis. And this being so, if the locus of $P$ be a given surface, then we
shall have a certain surface, the locus of $Q$; and so if the locus of $P$ be a given
curve in space, then we shall have a given curve in space, the locus of $Q$. In
particular, if the locus of $P$ be the plane of the three given points, then the locus
of $Q$ will be a certain surface, cutting the plane in a curve which is the locus in
the foregoing problem; and when $Q$ is situate on this curve, then also $P$ will be
situate on the same curve. Or if the locus of $P$ be the curve in question, then the
locus of $Q$ will be the same curve. Say, in general, that the loci of $P$ and $Q$ are
reciprocal loci, then the curve in the problem \textit{is its own reciprocal}. And we may
propose the following question:

Find the curve or surface, the locus of $P$, which is its own reciprocal.

We have also analogous to the original problem the following question in Solid
Geometry:

Given the four points $A$, $B$, $C$, $D$ in space, to find the locus of the points $P$, $Q$
such that
\[
PA + AQ = PB + BQ = PC + CQ = PD + DQ = \text{a given line.}\}
\]

\begin{center}
\textit{Solution by the Proposer.}
\end{center}

In general if $a$, $b$, $c$ be the sides of a triangle, and $f$, $g$, $h$ the distances of any
point from the angles of the triangle (or, what is the same thing, if $(a, b, c, f, g, h)$
be the distances of any four points in a plane from each other), then we have a certain
relation
\[
\phi\,(a, b, c, f, g, h) = 0.
\]
Hence if $r$, $s$, $t$ be the distances of the one focus from the angles of the triangle, and
the major axis is $=2\lambda$, then the distances for the other focus are $2\lambda-r$, $2\lambda-s$, $2\lambda-t$;
and considering the three angles and the other focus as a system of four points, we
have
\[
\phi\,(a, b, c, 2\lambda-r, 2\lambda-s, 2\lambda-t) = 0,
\]
which is a relation between the distances $r$, $s$, $t$ of the first focus from the angles of
the triangle, and which, treating these distances as coordinates (of course in a generalised
sense of the term ``Coordinate''), may be regarded as the equation of the required
locus. It is to be observed, that we have identically
\[
\phi\,(a, b, c, r, s, t) = 0,
\]
and the equation may be expressed in the simplified form
\[
\phi\,(a, b, c, 2\lambda-r, 2\lambda-s, 2\lambda-t) - \phi\,(a, b, c, r, s, t) = 0.
\]

To develope the solution, I notice that the expression for the equation $\phi\,(a, b, c, f, g, h)=0$ is
\begin{align*}
&b^2c^2\,(f^2+g^2) + c^2a^2\,(h^2+f^2) + a^2b^2\,(f^2+g^2)\\
&\quad{}+ g^2h^2\,(b^2+c^2) + h^2f^2\,(c^2+a^2) + f^2g^2\,(a^2+b^2)\\
&\quad{}- a^2f^2\,(a^2+f^2) - b^2g^2\,(b^2+g^2) - c^2h^2\,(c^2+h^2)\\
&\quad{}- a^2g^2h^2 - b^2h^2f^2 - c^2f^2g^2 - a^2b^2c^2 = 0;
\end{align*}
see my paper, ``Note on the value of certain determinants \&c.,'' \textit{Quart.\ Math.\ Journ.}\
vol.~\textsc{iii}. (1860), pp.\ 275---277, [286]. Or, as this may also be written
\[
\Sigma\,\{(b^2+c^2-a^2)\,(g^2h^2+a^2f^2) - a^2f^4\} - a^2b^2c^2 = 0,
\]
where $\Sigma$ refers to the simultaneous cyclical permutation of $(a, b, c)$ and of $(f, g, h)$.
Hence we have only in this equation to write $2\lambda-r$, $2\lambda-s$, $2\lambda-t$ in place of
$(f, g, h)$, and to omit the terms independent of $\lambda$, being in fact those which are equal
to $\phi\,(a, b, c, r, s, t)$. Observing that we have
\begin{align*}
g^2h^2+a^2f^2 &= [4\lambda^2-2\lambda(s+t)+st]^2 + a^2(2\lambda-r)^2\\
&= 16\lambda^4 - 16\lambda^3(s+t) + 4\lambda^2(s^2+t^2+4st+a^2) - 4\lambda\,[st(s+t)+a^2r] + s^2t^2+a^2r^2;\\
f^4 &= (2\lambda-r)^4 = 16\lambda^4 - 32\lambda^3r + 24\lambda^2r^2 - 8\lambda r^3 + r^4,
\end{align*}
the equation becomes
\begin{align*}
&16\lambda^4\,\{\Sigma\,(b^2+c^2-a^2) - \Sigma a^2\}\\
&{}- 16\lambda^3\,\{\Sigma\,(b^2+c^2-a^2)\,(s+t) - 2\Sigma a^2r\}\\
&{}+ 4\lambda^2\,\{\Sigma\,(b^2+c^2-a^2)\,(s^2+t^2+4st+a^2) - 6\Sigma a^2r^2\}\\
&{}- 4\lambda\,\{\Sigma\,(b^2+c^2-a^2)\,[st(s+t)+a^2r] - 2\Sigma a^2r^3\} = 0,
\end{align*}
where the $\Sigma$'s refer to the simultaneous cyclical permutation of the $(a, b, c)$ and the
$(r, s, t)$. The coefficients of $\lambda^4$ and $\lambda^3$ are, it is easy to see, each $=0$; and in the
coefficient of $\lambda^2$ the terms $\Sigma\,(b^2+c^2-a^2)\,(s^2+t^2) - 6\Sigma a^2r^2$ are $=-4\Sigma a^2r^2$; hence dividing
the whole equation by $4\lambda$, we find
\[
\lambda\,\{\Sigma\,(b^2+c^2-a^2)\,(4st+a^2) - 4\Sigma a^2r^2\} - \{\Sigma\,(b^2+c^2-a^2)\,[st(s+t)+a^2r] - 2\Sigma a^2r^3\} = 0,
\]
which is the required relation between $(r, s, t)$.

It may be noticed that, expressing the distances $r$, $s$, $t$ in terms of Cartesian or
trilinear coordinates $(x, y)$ or $(x, y, z)$, then $r^2$, $s^2$, $t^2$ are rational and integral functions
of the coordinates, and the form of the equation therefore is
\[
A_2 + B_2r + C_2s + D_2t + E_2st + F_2tr + G_2rs = 0,
\]
where the subscript numbers denote the degrees in regard to the coordinates. Multiply-
ing this equation successively by $1$, $r$, $s$, $t$, $st$, $tr$, $rs$, $rst$, we have eight equations linear
in the last-mentioned eight quantities, the coefficients being of known degrees respectively;
and eliminating the eight quantities, we have the rationalised equation expressed in the
form, determinant (of order $8$) $=0$; viz.\ this is
\[
\begin{vmatrix}
A_2, & B_2, & C_2, & D_2, & E_2, & F_2, & G_2, & 0 \\
B_2r^2, & A_2, & G_2r^2, & F_2r^2, & 0, & D_2, & C_2, & E_2 \\
C_2s^2, & G_2s^2, & A_2, & E_2s^2, & D_2, & 0, & B_2, & F_2 \\
D_2t^2, & F_2t^2, & E_2t^2, & A_2, & C_2, & B_2, & 0, & G_2 \\
E_2s^2t^2, & 0, & D_2t^2, & C_2s^2, & A_2, & G_2s^2, & F_2t^2, & B_2 \\
F_2t^2r^2, & D_2t^2, & 0, & B_2r^2, & G_2r^2, & A_2, & E_2t^2, & C_2 \\
G_2r^2s^2, & C_2s^2, & B_2r^2, & 0, & F_2r^2, & E_2s^2, & A_2, & D_2 \\
0, & E_2s^2t^2, & F_2t^2r^2, & G_2r^2s^2, & B_2r^2, & C_2s^2, & D_2t^2, & A_2
\end{vmatrix} = 0.
\]
This seems to be of the degree 18 in the coordinates, but it is probable that the real
degree is lower.

\subsection*{[Vol.~\textsc{iii}. pp.\ 63---65.]}

\paragraph*{1652.} (Proposed by W.~K.~Clifford.)---Through the angles $A$, $B$, $C$ of a plane
triangle three straight lines $Aa$, $Bb$, $Cc$ are drawn. A straight line $AR$ meets $Cc$ in
$R$; $RB$ meets $Aa$ in $P$; $PC$ meets $Bb$ in $Q$; $QA$ meets $Cc$ in $r$; and so on. Prove
that, after going twice round the triangle in this way, we always come back to the
same point.

Show that the theorem is its own reciprocal. Find the analogous properties of a
skew quadrilateral in space, and of a polygon of $n$ sides in a plane.

\begin{center}
\textit{Solution by Professor Cayley.}
\end{center}

1. The theorem may be thus stated: Given three lines $x$, $y$, $z$, and in these lines
respectively the points $A$, $B$, $C$; then there exist an infinity of hexagons, such that
the pairs of opposite angles lie in the lines $x$, $y$, $z$, respectively, and that the pairs
of opposite sides pass through the points $A$, $B$, $C$, respectively.

2. The demonstration is as follows: We have to show that, starting from an
arbitrary point $1$ in the line $x$, and constructing in the prescribed manner (as shown
successively in the figure) the points $2$, $3$, $4$, $5$, $6$, the last side $61$ of the hexagon
$123456$ will pass through $B$. By the construction, we have $A$, $2$, $3$ in a line, and
likewise $C$, $4$, $5$; hence, by Pascal's theorem, applied to the six points in a pair of
lines, the points of intersection of the lines $(25, 34)$, $(3C, A5)$, $(A4, C2)$, that is, the
points $B$, $6$, $1$, lie in a line; which is the required theorem.

3. More generally suppose that the points $A$, $B$, $C$ are not on the lines $x$, $y$, $z$,
respectively. I remark that it is not in general possible to describe a hexagon such
that the opposite angles lie in the lines $x$, $y$, $z$, respectively, and the opposite sides
pass through the points $A$, $B$, $C$, respectively; but if there exists one hexagon (viz.,
a proper hexagon, not a triangle twice repeated), then there exists an infinity of such
hexagons.

4. In fact, if it be required to find a polygon, the angles whereof lie in given
lines respectively, and the sides whereof pass through given points respectively; the
problem is either indeterminate or admits of only \textit{two} solutions. If therefore in any
particular case there are three or more solutions, the problem is indeterminate, and
has an infinity of solutions. Now, in the above-mentioned case of the three lines
and the three points, there exist \textit{two} triangles, the angles whereof lie in the given
lines, and the sides pass through the given points; and each triangle, taking the
angles twice over in the same order $123123$, is a hexagon satisfying the conditions of
the problem; hence, if we have besides a proper hexagon satisfying the conditions of
the problem, there are really \textit{three} solutions, and the problem is therefore indeterminate.

5. Suppose that the three lines $x$, $y$, $z$, and also two of the three points, say
the points $A$ and $B$, are given; we may construct geometrically a locus, such that,
taking for $C$ any point of this locus, the problem shall be indeterminate: in fact,
starting with the point $4$, and constructing successively the points $3$, $2$; taking an
arbitrary direction for the line $21$, and constructing successively the points $1$, $6$, $5$;
then the intersection of the lines $21$ and $54$ is a position of the point $C$: and by
taking any number of directions for the line $21$, we obtain for each of them a different
position of the point $C$; and so construct the locus.

6. The locus in question is, as will be shown, a line; and if the point $A$ is on
the line $x$, and the point $B$ on the line $y$, then the locus of $C$ will be the line $z$;
that is, $C$ being any point of the line $z$, the problem is indeterminate; which is
Mr Clifford's theorem.

7. To prove this, consider the lines $x$, $y$, $z$, and also the points $A$, $B$, $C$, as
given; the point $1$ is an arbitrary point on the line $x$, linearly determined by means
of a parameter $u$; and for every position of the point $1$ we have a corresponding
position of the point $4$; let $u'$ be the corresponding parameter for the point $4$; the
series of points $1$ is homographic with the series of points $4$; that is, the parameters
$u$, $u'$ are connected by an equation of the form $auu' + bu + cu' + d = 0$, (where of course
$a$, $b$, $c$, $d$ are functions of the parameters which determine the given lines $x$, $y$, $z$ and
the points $A$, $B$, $C$). But if the problem be indeterminate, then starting from the point
$1$ and constructing the point $4$, and again starting from the point $4$ and making the
very same construction, we arrive at the original point $1$, that is, $u$ must be the
same function of $u'$ that $u'$ is of $u$; and this will be the case if $b=c$; hence $b=c$
is the condition in order that the problem may be indeterminate.

8. To effect the calculation, take $x=0$, $y=0$, $z=0$ for the equations of the lines
$x$, $y$, $z$ respectively; and let $(\alpha, \beta, \gamma)$, $(\alpha', \beta', \gamma')$, $(\alpha'', \beta'', \gamma'')$ be the coordinates of the
points $A$, $B$, $C$ respectively. Let $1$ and $4$ be given as the intersections of the line
$x=0$ with the lines $y-uz=0$, $y-u'z=0$, respectively; and assume that for the point $2$
we have $y=0$, $z-vx=0$, and for the point $3$, $z=0$, $x-wy=0$. Then $1$, $C$, $2$ are in
a line; as are also $2$, $A$, $3$; $3$, $B$, $4$; hence we obtain
\[
v = \frac{\gamma''u - \beta''}{\alpha''u}, \quad w = \frac{\alpha v - \gamma}{\beta v}, \quad u' = \frac{\beta'w - \alpha}{\gamma'w};
\]
therefore, eliminating $v$ and $w$, we have
\[
(\alpha\gamma'' - \alpha''\gamma)\,\gamma'uu' - \alpha\beta''\gamma'u' - (\alpha\beta'\gamma'' - \alpha'\beta\gamma'' - \alpha''\beta'\gamma)\,u - \beta''(\alpha'\beta - \alpha\beta') = 0.
\]
The required condition, therefore, is
\[
\alpha\beta'\gamma'' = \alpha\beta'\gamma'' - \alpha'\beta\gamma'' - \alpha''\beta'\gamma, \quad \text{or}\quad \alpha\beta'\gamma'' - \alpha\beta''\gamma' - \alpha'\beta\gamma'' - \alpha''\beta'\gamma = 0;
\]
which is linear in regard to each of the three sets $(\alpha, \beta, \gamma)$, $(\alpha', \beta', \gamma')$, $(\alpha'', \beta'', \gamma'')$,
separately; that is, two of the points $A$, $B$, $C$ being given, the locus of the remaining
point is a line. In particular, if $\alpha=0$, $\beta'=0$; then the equation becomes $\alpha'\beta\gamma''=0$,
and assuming that neither $\alpha'=0$ or $\beta=0$, then the equation becomes $\gamma''=0$, that is,
$A$, $B$ being arbitrary points on the lines $x=0$, $y=0$ respectively, the locus of $C$ is the
line $z=0$.

9. Mr Clifford's theorem is clearly its own reciprocal. I do not know the \textit{precise}
analogues of his special form of the theorem; but the analogue of the more general
theorem stated in (6) is as follows: viz., we may have in the plane $n$ lines $x$, $y$, $z$, \ldots
and $n$ points $A$, $B$, $C$, \ldots, such that there exist an infinity of $2n$-gons whereof the
pairs of opposite angles lie in the given lines respectively; and the pairs of opposite
sides pass through the given points respectively; and if the $n$ lines and $n-1$ of the
$n$ points be assumed at pleasure, then the locus of the remaining point is a line. It
is moreover clear by the principle of reciprocity, that if the $n$ points and $n-1$ of
the $n$ lines be assumed at pleasure, then the envelope of the remaining line is a
point.

There exists also an analogue in space; viz.\ we may have $n$ lines $x$, $y$, $z$, \ldots and
$n$ lines $A$, $B$, $C$, \ldots such that there exist an infinity of (skew) $2n$-gons whereof the
pairs of opposite angles lie in the given lines $x$, $y$, $z$, \ldots respectively; and the pairs
of opposite sides meet in the given lines $A$, $B$, $C$, \ldots respectively. It may be added, that
if all but one of the $2n$ lines $x$, $y$, $z$, \ldots $A$, $B$, $C$ \ldots are given, then the `six coordi-
\noindent nates' of the remaining line satisfy a certain linear equation, but I do not stop to
explain the geometrical interpretation of this theorem.

\subsection*{[Vol.~\textsc{iii}. pp.\ 78, 79.]}

\paragraph*{1667.} (Proposed by Professor Sylvester.)

Show that the discriminant of the form
\[
ax^5 + b\lambda x^4y + c\lambda^2x^3y^2 + c\mu^2x^2y^3 + b\mu xy^4 + ay^5
\]
will be a rational integral function of the quantities $a$, $b$, $c$, $\lambda\mu$, $\lambda^5+\mu^5$, and of the
second degree only in respect to the last of them.

\begin{center}
\textit{Solution by Professor Cayley.}
\end{center}

In general
\[
\text{Disct.}\,(a, b, c, d, e, f\!\!\between\!\lambda x+\mu y,\ \lambda'x+\mu'y)^5 = (\lambda\mu' - \lambda'\mu)^{20}\,\text{Disct.}\,(a, b, c, d, e, f\!\!\between\! x, y)^5.
\]
Hence first, if $(\lambda, \mu, \lambda', \mu') = (0, 1, 1, 0)$, then
\[
\text{Disct.}\,(a, b, c, d, e, f\!\!\between\! y, x)^5 = \text{Disct.}\,(a, b, c, d, e, f\!\!\between\! x, y)^5;
\]
and secondly, if $\omega$ be an imaginary fifth root of unity and $(\lambda, \mu, \lambda', \mu') = (\omega, 0, 0, 1)$,
then
\[
\text{Disct.}\,(a, b, c, d, e, f\!\!\between\! \omega x, y)^5 = \text{Disct.}\,(a, b, c, d, e, f\!\!\between\! x, y)^5.
\]
These two results may also be written,
\begin{align*}
\text{Disct.}\,(a, b, c, d, e, f\!\!\between\! x, y)^5 &= \text{Disct.}\,(f, e, d, c, b, a\!\!\between\! x, y)^5,\\
\text{Disct.}\,(a, b, c, d, e, f\!\!\between\! x, y)^5 &= \text{Disct.}\,(a, b\omega^4, c\omega^3, d\omega^2, e\omega, f\!\!\between\! x, y)^5;
\end{align*}
that is, the discriminant of $(a, b, c, d, e, f\!\!\between\! x, y)^5$ is not altered by taking the coefficients in
a reverse order, or by multiplying the several coefficients by the powers $\omega^5$, $\omega^4$, $\omega^3$, $\omega^2$, $\omega$, of an
imaginary fifth root of unity. Applying these theorems to the form $(a, b\lambda, c\lambda^2, c\mu^2, b\mu, a\!\!\between\! x, y)^5$,
the discriminant is not altered by changing the coefficients into $(a, b\mu, c\mu^2, c\lambda^2, b\lambda, a)$;
that is, by the interchange of $\lambda$ and $\mu$; nor by changing the coefficients into
\[
(a, b\omega^4\lambda, c\omega^3\lambda^2, c\omega^2\mu^2, b\omega\mu, a), \quad \text{or}\quad \{a, b(\lambda\omega^4), c(\lambda\omega^4)^2, c(\mu\omega)^2, b(\mu\omega), a\}:
\]
that is, the discriminant is not altered by the change of $\lambda$, $\mu$ into $\lambda\omega^4$, $\mu\omega$ respectively. In
virtue of the second property the discriminant is a rational integral function of $(\lambda\mu, \lambda^5, \mu^5)$,
and then in virtue of the first property it is a rational integral function of $(\lambda\mu, \lambda^5\mu^5, \lambda^5+\mu^5)$,
that is, of $\lambda\mu$, $\lambda^5+\mu^5$. For the general form $(a, b, c, d, e, f\!\!\between\! x, y)^5$, if a term of the
discriminant be $a^\alpha b^\beta c^\gamma d^\delta e^\epsilon f^\phi$, then we have $\alpha+\beta+\gamma+\delta+\epsilon+\phi=8$, $5\alpha+4\beta+3\gamma+2\delta+\epsilon=20$;
hence attending only to the indices $\alpha$, $\beta$, $\gamma$ we have $5\alpha+4\beta+3\gamma > 20$, and therefore
\textit{\`{a} fortiori} $3\beta+3\gamma > 20$, so that $\beta+\gamma$ is $=6$ at most. It follows that for the form
$(a, b\lambda, c\lambda^2, c\mu^2, b\mu, a\!\!\between\! x, y)^5$, the sum of the indices of $b\lambda$, $c\lambda^2$ is $=6$ at most, and
therefore, even if the index of $c\lambda^2$ is $=6$, the index of $\lambda$ will be only $=12$, that is, the
discriminant contains no power of $\lambda$ higher than $\lambda^{12}$: hence considered as a function of
$\lambda\mu$, $\lambda^5+\mu^5$, the highest power of $\lambda^5+\mu^5$ is $(\lambda^5+\mu^5)^2$; which completes the theorem.

\subsection*{[Vol.~\textsc{iii}. p.\ 90.]}

\paragraph*{1687.} (Proposed by Professor Cayley.)---To describe a spherical triangle such that
the angles thereof and of the polar triangle lie on a spherical conic.

On the sphere, the locus of a point such that the perpendiculars from it upon
the sides of a given spherical triangle have their feet on a line (great circle), is in
general a spherical cubic; if however the triangle be such as is mentioned in the
above Problem, then the locus breaks up into a line (great circle) and into the conic
through the angles of the given and polar triangles.

\subsection*{[Vol.~\textsc{iii}. pp.\ 92---96.]}

\paragraph*{1690.} (Proposed by W.~A.~Whitworth, M.A.)---If $ABC$ be the triangle formed by
the three diagonals $aa'$, $bb'$, $cc'$ of a complete quadrilateral $aa'bb'cc'$, then a conic can
be found having double contact in the chord $aa'$ with the critical conic of the quadri-
lateral $bb'cc'$, double contact in the chord $bb'$ with the critical conic of the quadrilateral
$cc'aa'$, and double contact in the chord $cc'$ with the critical conic of the quadrilateral
$aa'bb'$.

\begin{center}
\textit{3. Solution by Professor Cayley.}
\end{center}

1. The equations of the sides of the quadrilateral may be taken to be respectively
$x=0$, $y=0$, $z=0$, $w=0$, where the implicit constants are so determined that we have
identically
\[
x + y + z + w = 0;
\]
this being so, the equations of the three diagonals are respectively
\begin{alignat*}{3}
x+y&=0, &\quad \text{or}\quad z+w&=0, &\quad \text{or}\quad x+y-z-w&=0 \quad \text{(three equivalent forms)}\\
x+z&=0, &\quad \text{or}\quad y+w&=0, &\quad \text{or}\quad x-y+z-w&=0 \quad \text{(}\quad\;\text{``}\quad\quad\;\text{``}\quad\quad\;\text{)}\\
x+w&=0, &\quad \text{or}\quad y+z&=0, &\quad \text{or}\quad x-y-z+w&=0 \quad \text{(}\quad\;\text{``}\quad\quad\;\text{``}\quad\quad\;\text{)}
\end{alignat*}
and the equations of the critical conics are respectively
\[
xy + zw = 0, \quad xz + yw = 0, \quad xw + yz = 0.
\]
Hence we see that the equation of the required conic is
\[
\Delta = x^2 + y^2 + z^2 + w^2 - 2yz - 2zx - 2xy - 2xw - 2yw - 2zw = 0.
\]
In fact this equation may be written
\begin{align*}
\Delta &= (x+y-z-w)^2 - 4(xy+zw) = 0,\\
\Delta &= (x-y+z-w)^2 - 4(xz+yw) = 0,\\
\Delta &= (x-y-z+w)^2 - 4(xw+yz) = 0,
\end{align*}
equations which put in evidence the double contact with the three critical conics
respectively. We have also, identically,
\[
\Delta = (x+y+z+w)(x+y-3z-w) - 2w(x+y-z-w) + 4(z^2-xy),
\]
and the equation $\Delta=0$ may therefore be written
\[
\Delta = -2w(x+y-z-w) + 4(z^2-xy) = 0,
\]
a form which shows that the conic $z^2-xy=0$ meets the line $w=0$ in the same two
points in which it is met by the conic $\Delta=0$. And it hence appears by symmetry that
the conics
\begin{align*}
\Delta=0,\ x^2-yz=0,\ y^2-zx=0,\ z^2-xy=0 &\text{ meet the line } w=0 \text{ in the same two points},\\
\Delta=0,\ w^2-yz=0,\ y^2-zw=0,\ z^2-wy=0 &\text{ meet the line } x=0 \text{ in the same two points},\\
\Delta=0,\ w^2-xz=0,\ x^2-zw=0,\ z^2-wx=0 &\text{ meet the line } y=0 \text{ in the same two points},\\
\Delta=0,\ w^2-xy=0,\ x^2-yw=0,\ y^2-wx=0 &\text{ meet the line } z=0 \text{ in the same two points},
\end{align*}
which are the relations constituting the latter part of the proposed theorem.

2. The analogous theorems in space may be briefly referred to. Taking $x=0$,
$y=0$, $z=0$, $w=0$ as the equations of the faces of a tetrahedron $ABCD$, then the
implicit constants may be so determined that the coordinates of a given arbitrary point
$O$ shall be $(1, 1, 1, 1)$. We may by lines drawn from the vertices of the tetrahedron
project the point $O$ on the faces, so as to obtain a point in each of the four faces;
and then in each face, by lines drawn from the vertices of the face, project the point
in that face upon the edges of the face; the two points thus obtained on each edge
of the tetrahedron are (it is easy to see) one and the same point; that is, we have
on each edge of the tetrahedron a point; and there exists a quadric surface
\[
\Delta = x^2 + y^2 + z^2 + w^2 - 2yz - 2zx - 2xy - 2xw - 2yw - 2zw = 0
\]
touching the edges of the tetrahedron in these six points respectively.

The surface in question has plane contact with
\begin{alignat*}{5}
&\text{the hyperboloid } &&xy+zw=0 &&\text{along the intersection with } &&x+y-z-w=0,\\
&\text{``}\qquad\;\;\; &&xz+yw=0 &&\text{``}\qquad\qquad\;\;\;\text{``} &&x-y+z-w=0,\\
&\text{``}\qquad\;\;\; &&xw+yz=0 &&\text{``}\qquad\qquad\;\;\;\text{``} &&x-y-z+w=0,
\end{alignat*}
and moreover the surfaces
\begin{align*}
\Delta=0,\ x^2-yz=0,\ y^2-zx=0,\ z^2-xy=0 &\text{ meet the line } w=0,\ x+y+z+w=0\\
&\qquad\qquad\qquad\qquad\qquad\text{in the same two points};\\
\Delta=0,\ w^2-yz=0,\ y^2-zw=0,\ z^2-wy=0 &\text{ meet the line } x=0,\ x+y+z+w=0\\
&\qquad\qquad\qquad\qquad\qquad\text{in the same two points};\\
\Delta=0,\ w^2-xz=0,\ x^2-zw=0,\ z^2-wx=0 &\text{ meet the line } y=0,\ x+y+z+w=0\\
&\qquad\qquad\qquad\qquad\qquad\text{in the same two points};\\
\Delta=0,\ w^2-xy=0,\ x^2-yw=0,\ y^2-wx=0 &\text{ meet the line } z=0,\ x+y+z+w=0\\
&\qquad\qquad\qquad\qquad\qquad\text{in the same two points}.
\end{align*}

With respect to the construction of the four planes,
\[
x+y-z-w=0,\quad x-y+z-w=0,\quad x-y-z+w=0,\quad x+y+z+w=0,
\]
it is to be observed that if through any edge of the tetrahedron, for instance the
edge $x=0$, $y=0$, we draw the plane $x-y=0$ through the point $O$, then the harmonic
of this in regard to the planes $x=0$, $y=0$ is the plane $x+y=0$; we have thus six
planes, one through each edge of the tetrahedron, viz., these are $y+z=0$, $z+x=0$,
$x+y=0$, $x+w=0$, $y+w=0$, $z+w=0$; the six planes being the faces of a hexahedron,
which is such that the vertices of the tetrahedron are four of the eight vertices of
the hexahedron. The pairs of opposite faces of the hexahedron meet in three lines
lying in the plane $x+y+z+w=0$, and consequently forming a triangle such that
through each side of the triangle there pass two opposite faces of the hexahedron;
the planes $x+y-z-w=0$, $x-y+z-w=0$, $x-y-z+w=0$ are the harmonics of the
plane $x+y+z+w=0$ in respect of the pairs of opposite faces of the hexahedron;
viz., the plane $x+y-z-w=0$ is the harmonic of the plane $x+y+z+w=0$ in respect
to the planes $x+y=0$, $z+w=0$; and the like for the other two planes $x-y+z-w=0$
and $x-y-z+w=0$ respectively.

\subsection*{[Vol.~\textsc{iv}. July to December, 1865, pp.\ 17, 18.]}

\paragraph*{1710.} (Proposed by Professor Cayley.)---Trace the curve $y^4 - 2y^2zx - z^4 = 0$, where
the coordinates are such that $x+y+z=0$ is the line \textit{infinity}.

\begin{center}
\textit{Solution by the Proposer.}
\end{center}

We have $x = \dfrac{y^4-z^4}{2y^2z}$; or writing $y=\theta z$, then $x = \dfrac{\theta^4-1}{2\theta^2}\,z$, that is
\[
x : y : z = \theta^4-1 : 2\theta^3 : 2\theta^2.
\]
Hence, we see that $y$, $z$ are indefinitely small in comparison of $x$,
\begin{alignat*}{2}
&\text{if } \theta=\infty, \text{ and then } &&x:y:z = \theta^4 : 2\theta^3 : 2\theta^2, \text{ that is } y^2=2zx;\\
&\text{or, if } \theta=0, \text{ and then } &&x:y:z = -1 : 2\theta^3 : 2\theta^2, \text{ that is } z^3 = -2y^2x;
\end{alignat*}
so that in the neighbourhood of the point $(y=0, z=0)$ there are two branches coin-
\noindent ciding with the parabola $y^2=2zx$ and with the semicubical parabola $z^3=-2y^2x$,
respectively.

To find the points at infinity we have $x+y+z=0$, that is $\theta^4 + 2\theta^3 + 2\theta^2 - 1 = 0$
$= (\theta+1)(\theta^3 + \theta^2 + \theta - 1) = 0$; and observing that the equation $\theta^3 + \theta^2 + \theta - 1 = 0$ has one
real root, say $\theta=k$, if $k$ be the real root of the equation $k^3+k^2+k-1=0$ ($k=.505$
nearly),---there are two real points at infinity, viz., corresponding to $\theta=-1$, we have
the point $(0, -1, 1)$, and corresponding to $\theta=k$ the point $(-1-k, k, 1)$.

The equation of the tangent at a point $(\alpha, \beta, \gamma)$ is
\[
x\,(-\beta^2\gamma) + y\,(2\beta^3 - 2\alpha\beta\gamma) + z\,(-\alpha\beta^2 - 2\gamma^3) = 0,
\]
and hence writing $(\alpha, \beta, \gamma) = (0, -1, 1)$ we have the asymptote $x+2y+2z=0$: to find
where this meets the curve, we have $\theta^4 + 4\theta^3 + 4\theta^2 - 1 = 0$, that is $(\theta+1)^2(\theta^2+2\theta-1)=0$,
or at the points of intersection $\theta^2+2\theta-1=0$, that is $\theta=-1\pm\surd2$, or there are two
real points of intersection.

Again writing $(\alpha, \beta, \gamma) = (-1-k, k, 1)$ we find the asymptote $k^2x - 2y + (k+1)z = 0$:
to find where this meets the curve, we have $k^2(\theta^4-1) - 4k\theta^3 + (2k+2)\,\theta^2 = 0$, that is
$k^2\theta^4 - 4k\theta^3 + (2k+2)\,\theta^2 - k^2 = (\theta-k)^2\,[k^2\theta^2 - 2(k^2+k+1)\,\theta - 1] = 0$; or for the points of
intersection $k^2\theta^2 - 2(k^2+k+1)\,\theta - 1 = 0$, an equation in $\theta$ with two real roots, hence
the points of intersection are real.

It is now easy to lay down the curve; viz., if, to fix the ideas, the fundamental
triangle is taken to be equilateral, and the coordinates $x$, $y$, $z$ are considered to be
positive for points \textit{within} the triangle, then the curve is as shown in the annexed
figure~1.

\begin{center}
\textit{[Figure 1: The curve $y^4 - 2y^2zx - z^4 = 0$ drawn within an equilateral triangle,
showing two infinite branches passing through the triangle with asymptotes.
The curve has a branch near the vertex $x$ coinciding with the parabola $y^2=2zx$
and another branch coinciding with the semicubical parabola $z^3=-2y^2x$.]}
\end{center}

It may be remarked that the curve is met by every real line in two real points
at least, and consequently that it is not the projection of any finite curve whatever.
By a modification of the constants of the equation, we might obtain curves which are
finite, such as the curve in figure~2; or curves with two or four infinite branches,
which are the projections of such a finite curve.

\begin{center}
\textit{[Figure 2: A finite modification of the curve, forming a closed oval shape.]}
\end{center}

\subsection*{[Vol.~\textsc{iv}. pp.\ 32---37.]}

\paragraph*{1744.} (Proposed by W.~S.~Burnside, B.A.)---It is required to find $(x_1, y_1, z_1)$,
functions of $(x, y, z)$, such that we may have identically
\[
\frac{x_1{}^3 + y_1{}^3 + z_1{}^3}{x_1y_1z_1} = \frac{x^3 + y^3 + z^3}{xyz}.
\]

\begin{center}
\textit{Solution by Professor Cayley.}
\end{center}

The Solution is in fact given in my ``Memoir on Curves of the Third Order,''
\textit{Philosophical Transactions}, vol.~cxlvii. (1857), pp.\ 415---446, [146].

Write $\dfrac{x^3+y^3+z^3}{xyz} = -6l$; then, taking $(X, Y, Z)$ as current coordinates, $(x, y, z)$ are,
it is clear, the coordinates of a point on the cubic curve $X^3+Y^3+Z^3+6lXYZ=0$;
and if $(x_1, y_1, z_1)$ are the coordinates of any other point on the same cubic curve, then
we shall have
\[
\frac{x_1{}^3+y_1{}^3+z_1{}^3}{x_1y_1z_1} = -6l = \frac{x^3+y^3+z^3}{xyz},
\]
so that $(x_1, y_1, z_1)$ will satisfy the condition in question. Hence, if from a given point
$(x, y, z)$ on the cubic curve we obtain by any geometrical construction another point
on the curve, the coordinates of this new point will be functions (and, if the construction
is such as to lead to a single point only, they will be \textit{rational} functions) of $(x, y, z)$,
satisfying the condition in question.

For instance, if the point $(x, y, z)$ be joined with any point $(\alpha, \beta, \gamma)$ on the curve,
the joining line will again meet the curve in a single point, which may be taken to
be the point $(x_1, y_1, z_1)$. This assumes that we know on the cubic curve a point
$(\alpha, \beta, \gamma)$; but such a point at once presents itself, viz., we may write $(\alpha, \beta, \gamma) = (1, -1, 0)$;
which gives only the self-evident solution $(x_1, y_1, z_1) = (y, x, z)$. The point $(1, -1, 0)$
is clearly one of the nine points of inflexion of the cubic curve, and by using these
in any manner whatever, viz., joining the point $(x, y, z)$ with any point of inflexion,
and then the new point with any other point of inflexion, and so on indefinitely, we
obtain in connexion with the given point $(x, y, z)$ seventeen other points on the curve,
in all a system of eighteen points: these are
\begin{alignat*}{4}
&(x, y, z), \quad &&(x, \omega y, \omega^2z), \quad &&(x, \omega^2y, \omega z), \quad &&(x, z, y), \quad (x, \omega z, \omega^2y), \quad (x, \omega^2z, \omega y)\\
&(y, z, x), \quad &&(\omega y, \omega^2z, x), \quad &&(\omega^2y, \omega z, x), \quad &&(z, y, x), \quad (\omega z, \omega^2y, x), \quad (\omega^2z, \omega y, x)\\
&(z, x, y), \quad &&(\omega^2z, x, \omega y), \quad &&(\omega z, x, \omega^2y), \quad &&(y, x, z), \quad (\omega^2y, x, \omega z), \quad (\omega y, x, \omega^2z)
\end{alignat*}
possessing remarkable geometrical properties; and of course each of the seventeen new
points furnishes a (self-evident) solution of the given identity.

But we may take $(\alpha, \beta, \gamma) = (x, y, z)$; the point $(x_1, y_1, z_1)$ is here the point of
intersection of the cubic by the tangent at the point $(x, y, z)$; or say it is the
``tangential'' of the point $(x, y, z)$. The values thus obtained for $(x_1, y_1, z_1)$ are
\[
(x_1, y_1, z_1) = \{x(y^3-z^3),\ y(z^3-x^3),\ z(x^3-y^3)\},
\]
which (excluding the above-mentioned self-evident solutions) is in fact the most simple
solution of the proposed identity. In order to verify that the last-mentioned values
of $(x_1, y_1, z_1)$ are in fact the coordinates of the tangential of $(x, y, z)$, I observe that
this will be the case if only we have
\[
(x^3+2lxyz)\,x_1 + (y^3+2lxyz)\,y_1 + (z^3+2lxyz)\,z_1 = 0, \quad x_1{}^3+y_1{}^3+z_1{}^3+6lx_1y_1z_1=0,
\]
the first of which is obviously satisfied by the values in question; and for the verification
of the second equation,
\begin{align*}
x_1{}^3+y_1{}^3+z_1{}^3 &= x^3(y^3-z^3)^3 + y^3(z^3-x^3)^3 + z^3(x^3-y^3)^3,\\
&= -x^3(y^3-z^3) - y^3(z^3-x^3) - z^3(x^3-y^3),\\
&= (x^3+y^3+z^3)\,(y^3-z^3)\,(z^3-x^3)\,(x^3-y^3),\\
x_1y_1z_1 &= xyz\,(y^3-z^3)\,(z^3-x^3)\,(x^3-y^3),
\end{align*}
therefore
\[
x_1{}^3+y_1{}^3+z_1{}^3+6lx_1y_1z_1 = (x^3+y^3+z^3+6lxyz)\,(y^3-z^3)\,(z^3-x^3)\,(x^3-y^3) = 0
\]
if $x^3+y^3+z^3+6lxyz=0$; the same equations verify at once the identity
\[
\frac{x_1{}^3+y_1{}^3+z_1{}^3}{x_1y_1z_1} = \frac{x^3+y^3+z^3}{xyz}.
\]

Another solution is as follows: viz., if we take the third intersection with the
cubic of the line joining the points $(y, x, z)$ and $\{x(y^3-z^3),\, y(z^3-x^3),\, z(x^3-y^3)\}$, the
coordinates of the line in question are
\begin{alignat*}{2}
x_1 : y_1 : z_1 &=\; &&x^6y^3 + y^6z^3 + z^6x^3 - 3x^3y^3z^3\\
&:\; &&x^3y^6 + y^3z^6 + z^3x^6 - 3x^3y^3z^3\\
&:\; &&xyz\,(x^6+y^6+z^6 - y^3z^3 - z^3x^3 - x^3y^3).
\end{alignat*}
According to a very beautiful theorem of Professor Sylvester's in relation to the theory of
cubic curves, the coordinates of a point which depends linearly on a given point of the
curve are necessarily rational and integral functions of a \textit{square degree} of the coordi-
\noindent nates $(x, y, z)$ of the given point; and moreover that (considering as \textit{one} solution those
which can be derived from each other by a mere permutation of the coordinates, or
change of $x$ into $\omega x$, \&c.), there is only one solution of a given square degree $m^2$; the
solutions of the degrees $4$ and $9$ are given above. The tangential of the tangential, or
second tangential of the point $(x, y, z)$, gives the solution of the degree $16$; joining
this second tangential with the original point $(x, y, z)$, we have the solution of the
degree $25$; and the same solution is also given as the sixth point of intersection with
the cubic, of the conic of $5$-pointic intersection at the point $(x, y, z)$. See my memoir
``On the conic of $5$-pointic contact at any point of a plane curve,'' \textit{Phil.\ Trans.}\ vol.\ cxlx. (1859), pp.\ 371---400, [261].

\vspace{1em}
\begin{center}
\textit{Addition to the foregoing Solution. On a system of Eighteen Points on a Cubic Curve.}
\end{center}
\vspace{0.5em}

Considering the cubic curve $x^3+y^3+z^3+6lxyz=0$, we have the nine points of
inflexion, which I represent as follows:
\begin{alignat*}{3}
a &= (0, 1, -1), &\quad d &= (-1, 0, 1), &\quad g &= (1, -1, 0),\\
b &= (0, 1, -\omega), &\quad e &= (-\omega, 0, 1), &\quad h &= (1, -\omega, 0),\\
c &= (0, 1, -\omega^2), &\quad f &= (-\omega^2, 0, 1), &\quad i &= (1, -\omega^2, 0),
\end{alignat*}
viz., $\omega$ being an imaginary cube root of unity, the coordinates of $a$ are $(0, 1, -1)$, those
of $b$, $(0, 1, -\omega)$, \&c.

The points of inflexion lie (as is known) by threes on twelve lines; viz., the lines are
\begin{alignat*}{4}
&abc, &\quad &afh, &\quad &bfg, &\quad &cei,\\
&adg, &\quad &bdi, &\quad &cdh, &\quad &def,\\
&aei, &\quad &beh, &\quad &ceg, &\quad &ghi.
\end{alignat*}

Consider now a point on the curve, the coordinates whereof are $(x, y, z)$, where of
course $x^3+y^3+z^3+6lxyz=0$; this is one of a system of eighteen points on the curve,
which may be represented as follows:
\begin{alignat*}{3}
A &= (x, y, z), &\quad D &= (x, \omega y, \omega^2z), &\quad G &= (x, \omega^2y, \omega z),\\
B &= (y, z, x), &\quad E &= (\omega y, \omega^2z, x), &\quad H &= (\omega^2y, \omega z, x),\\
C &= (z, x, y), &\quad F &= (\omega^2z, x, \omega y), &\quad I &= (\omega z, x, \omega^2y),\\
J &= (x, z, y), &\quad M &= (x, \omega z, \omega^2y), &\quad P &= (x, \omega^2z, \omega y),\\
K &= (z, y, x), &\quad N &= (\omega z, \omega^2y, x), &\quad Q &= (\omega^2z, \omega y, x),\\
L &= (y, x, z), &\quad O &= (\omega^2y, x, \omega z), &\quad R &= (\omega y, x, \omega^2z),
\end{alignat*}
viz., the coordinates of $A$ are $(x, y, z)$; those of $B$ are $(y, z, x)$, \&c.

The tangent at $A$ meets the curve in a point, ``the tangential of $A$,'' the coordi-
\noindent nates whereof are $x(y^3-z^3)$, $y(z^3-x^3)$, $z(x^3-y^3)$; which point may be called $A'$. And
we have thus the eighteen tangentials
\[
A',\ B',\ C',\ D',\ E',\ F',\ G',\ H',\ I',\ J',\ K',\ L',\ M',\ N',\ O',\ P',\ Q',\ R'.
\]
The eighteen points $A$, $B$, \&c., have the following property; viz., the line joining any
two of them meets the cubic in a third point, which is either one of the nine points
of inflexion, or one of the eighteen tangentials; there are through each point of inflexion
$9$ such lines, and through each tangential $4$ such lines; $(9 \times 9) + (18 \times 4) = 153 = \frac12(18.17)$,
the number of pairs of points $AB$, $AC$, \&c. The lines through the inflexions are the
$81$ lines obtained by joining any one of the points $(A, B, C, D, E, F, G, H, I)$ with
any one of the points $(J, K, L, M, N, O, P, Q, R)$, as shown in the following Table:

\begin{center}
\begin{tabular}{c|ccccccccc}
 & $A$ & $B$ & $C$ & $D$ & $E$ & $F$ & $G$ & $H$ & $I$ \\
\hline
$J$ & $a$ & $d$ & $g$ & $c$ & $f$ & $i$ & $b$ & $e$ & $h$ \\
$K$ & $d$ & $g$ & $a$ & $f$ & $i$ & $c$ & $e$ & $h$ & $b$ \\
$L$ & $g$ & $a$ & $d$ & $i$ & $c$ & $f$ & $h$ & $b$ & $e$ \\
$M$ & $c$ & $f$ & $i$ & $b$ & $e$ & $h$ & $a$ & $d$ & $g$ \\
$N$ & $f$ & $i$ & $c$ & $e$ & $h$ & $b$ & $d$ & $g$ & $a$ \\
$O$ & $i$ & $c$ & $f$ & $h$ & $b$ & $e$ & $g$ & $a$ & $d$ \\
$P$ & $b$ & $e$ & $h$ & $a$ & $d$ & $g$ & $c$ & $f$ & $i$ \\
$Q$ & $e$ & $h$ & $b$ & $d$ & $g$ & $a$ & $f$ & $i$ & $c$ \\
$R$ & $h$ & $b$ & $e$ & $g$ & $a$ & $d$ & $i$ & $c$ & $f$
\end{tabular}
\end{center}

\end{document}
